\section{AMC 10 10B 2018}
        \begin{enumerate}
        	\item (\textit{AMC 10 10B 2018}) Kate bakes a $20$-inch by $18$-inch pan of cornbread. The cornbread is cut into pieces that measure $2$ inches by $2$ inches. How many pieces of cornbread does the pan contain?


 $\textbf{(A) } 90 \qquad \textbf{(B) } 100 \qquad \textbf{(C) } 180 \qquad \textbf{(D) } 200 \qquad \textbf{(E) } 360$

 \textbf{Answer: }A



\textbf{Solution 1}

The area of the pan is $20\cdot18=360$. Since the area of each piece is $2\cdot2=4$, there are $\frac{360}{4} = \boxed{\textbf{(A) } 90}$ pieces.

 



\textbf{Solution 2}

By dividing each of the dimensions by $2$, we get a $10\times9$ grid that makes $\boxed{\textbf{(A) } 90}$ pieces.

 


	\item (\textit{AMC 10 10B 2018}) Sam drove $96$ miles in $90$ minutes. His average speed during the first $30$ minutes was $60$ mph (miles per hour), and his average speed during the second $30$ minutes was $65$ mph. What was his average speed, in mph, during the last $30$ minutes?


 $\textbf{(A) } 64 \qquad \textbf{(B) } 65 \qquad \textbf{(C) } 66 \qquad \textbf{(D) } 67 \qquad \textbf{(E) } 68$

 \textbf{Answer: }D



\textbf{Solution 1}

Suppose that Sam's average speed during the last $30$ minutes was $x$ mph.

 Recall that a half hour is equal to $30$ minutes. Therefore, Sam drove $60\cdot0.5=30$ miles during the first half hour, $65\cdot0.5=32.5$ miles during the second half hour, and $x\cdot0.5$ miles during the last half hour. We have 
 \begin{align*} 30+32.5+x\cdot0.5&=96 \\ x\cdot0.5&=33.5 \\ x&=\boxed{\textbf{(D) } 67}. \end{align*}
~Haha0201 ~MRENTHUSIASM

 



\textbf{Solution 2}

Suppose that Sam's average speed during the last $30$ minutes was $x$ mph.

 Note that Sam's average speed during the entire trip was $\frac{96}{3/2}=64$ mph. Since Sam drove at $60$ mph, $65$ mph, and $x$ mph for the same duration ($30$ minutes each), his average speed during the entire trip was the average of $60$ mph, $65$ mph, and $x$ mph. We have
 \begin{align*} \frac{60+65+x}{3}&=64 \\ 60+65+x&=192 \\ x&=\boxed{\textbf{(D) } 67}. \end{align*}
~coolmath\_2018 ~MRENTHUSIASM

 


	\item (\textit{AMC 10 10B 2018}) In the expression $\left(\underline{\qquad}\times\underline{\qquad}\right)+\left(\underline{\qquad}\times\underline{\qquad}\right)$ each blank is to be filled in with one of the digits $1,2,3,$ or $4,$ with each digit being used once. How many different values can be obtained?


 $\textbf{(A) }2 \qquad \textbf{(B) }3 \qquad \textbf{(C) }4 \qquad \textbf{(D) }6 \qquad \textbf{(E) }24 \qquad$

 \textbf{Answer: }B



\textbf{Solution 1 - Combinations}

We have $\binom{4}{2}$ ways to choose the pairs, and we have $2!$ ways for the values to be rearranged, hence $\frac{6}{2}=\boxed{\textbf{(B) }3}$




\textbf{Solution 2}

We have four available numbers $(1, 2, 3, 4)$. Because different permutations do not matter because they are all addition and multiplication, if we put $1$ on the first space, it is obvious there are $\boxed{\textbf{(B) }3}$ possible outcomes $(2, 3, 4)$.

 



\textbf{Solution 3}

There are exactly $4!$ ways to arrange the numbers and $2!2!2!$ overcounts per way due to commutativity. Therefore, the answer is $\frac{4!}{2!2!2!}=\boxed{\textbf{(B) }3}$



	\item (\textit{AMC 10 10B 2018}) A three-dimensional rectangular box with dimensions $X$, $Y$, and $Z$ has faces whose surface areas are $24, 24, 48, 48, 72,$ and $72$ square units. What is $X+Y+Z$?


 $\textbf{(A) }18 \qquad \textbf{(B) }22 \qquad \textbf{(C) }24 \qquad \textbf{(D) }30 \qquad \textbf{(E) }36 \qquad$

 \textbf{Answer: }B



\textbf{Solution 1}

Let $X$ be the length of the shortest dimension and $Z$ be the length of the longest dimension. Thus, $XY = 24$, $YZ = 72$, and $XZ = 48$.Divide the first two equations to get $\frac{Z}{X} = 3$. Then, multiply by the last equation to get $Z^2 = 144$, giving $Z = 12$. Following, $X = 4$ and $Y = 6$.

 The final answer is $4 + 6 + 12 = 22$. $\boxed{B}$




\textbf{Solution 2}

Simply guess and check to find that the dimensions are $4$ by $6$ by $12$. Therefore, the answer is $4 + 6 + 12 = 22$. $\boxed{B}$




\textbf{Solution 3}

If you find the GCD of $24$, $48$, and $72$ you get your first number, $12$. After this, do $48 \div 12$ and $72 \div 12$ to get $4$ and $6$, the other 2 numbers. When you add up your $3$ numbers, you get $22$ which is $\boxed{B}$.

 



\textbf{Solution 4}

Since the surface areas of the faces are the product of two of the dimensions. Therefore, $XY=24$, $XZ=48$, and $YZ=72$. You can multiply $XY \times XZ \times YZ$, which simplifies to ${XYZ}^2=24 \times 48 \times 72$ which means that the volume$XYZ$ equals$\sqrt{24*48*72}=\sqrt{24^2*12^2}=24*12=288$. The individual dimensions, $X$, $Y$, and $Z$ can be found by doing $\frac{XYZ}{XY}$, $\frac{XYZ}{YZ}$, and $\frac{XYZ}{XZ}$, which yields $Z=12$, $Y=6$, and $X=4$. Adding this up, we have that $X+Y+Z=22$ which is $\boxed{B}$- Solution by smartninja2000

 


	\item (\textit{AMC 10 10B 2018}) How many subsets of $\{2,3,4,5,6,7,8,9\}$ contain at least one prime number?


 $\textbf{(A) }128 \qquad \textbf{(B) }192 \qquad \textbf{(C) }224 \qquad \textbf{(D) }240 \qquad \textbf{(E) }256 \qquad$

 \textbf{Answer: }D



\textbf{Solution 1 - Complementary Counting}

We use complementary counting, or 

 $\text{total - what we don't want = what we want}$. 

 There are a total of $2^8$ ways to create subsets (consider including or excluding each number) and there are a total of $2^4$ subsets only containing composite numbers (the composite numbers are $4, 6, 8, 9$). Therefore, there are $2^8-2^4=\boxed{\textbf{(D) }240}$ total ways to have at least one prime in a subset.

 



\textbf{Solution 2 (Using Answer Choices)}

Well, there are 4 composite numbers, and you can list them in a 1 number format, a 2 number, 3 number, and a 4 number format. Now, we can use combinations.

 $\binom{4}{1} + \binom{4}{2} + \binom{4}{3} + \binom{4}{4} = 15$. Using the answer choices, the only multiple of 15 is $\boxed{\textbf{(D) }240}$

 By: K6511

 



\textbf{Solution 3}

Subsets of $\{2,3,4,5,6,7,8,9\}$ include a single digit up to all eight numbers. Therefore, we must add the combinations of all possible subsets and subtract from each of the subsets formed by the composite numbers.

 Hence:

 $\binom{8}{1} - \binom{4}{1} + \binom{8}{2} - \binom{4}{2} + \binom{8}{3} - \binom{4}{3} + \binom{8}{4} - 1 + \binom{8}{5} + \binom{8}{6} + \binom{8}{7} + 1 = \boxed{\textbf{(D) }240}$

 By: pradyrajasai

 


 



\textbf{Solution 4}

Total subsets is $(2^8) = 256$ Using complementary counting and finding the sets with composite numbers:only 4,6,8 and 9 are composite. Each one can be either in the set or out:$2^4$ = 16$256-16=240$$\boxed{\textbf{(D) }240}$

 -goldenn

 



\textbf{Solution 5}

We multiply the number of possibilities of the set having prime numbers and the set having composites.

 The possibilities of primes are $2^4-1=15$ (As there is one solution not containing any primes)

 The possibilities of the set containing composites are $2^4=16$ (There can be a set with no composites)

 Multiplying this we get $15 \cdot 16 = \boxed{\textbf{(D) }240}$. 

 -middletonkids

 


	\item (\textit{AMC 10 10B 2018}) A box contains $5$ chips, numbered $1, 2, 3, 4,$ and $5$. Chips are drawn randomly one at a time without replacement until the sum of the values drawn exceeds $4$. What is the probability that $3$ draws are required?


 $\textbf{(A) }\frac{1}{15} \qquad \textbf{(B) }\frac{1}{10} \qquad \textbf{(C) }\frac{1}{6} \qquad \textbf{(D) }\frac{1}{5} \qquad \textbf{(E) }\frac{1}{4} \qquad$

 \textbf{Answer: }D



\textbf{Solution 1}

Notice that the only four ways such that $3$ draws are required are $1,2$; $1,3$; $2,1$; and $3,1$. Notice that each of those cases has a $\frac{1}{5} \cdot \frac{1}{4}$ chance, so the answer is $\frac{1}{5} \cdot \frac{1}{4} \cdot 4 = \frac{1}{5}$, or $\boxed{D}$.

 Jonathan Xu (pi\_is\_delicious\_69420)

 



\textbf{Solution 1.1}

Since our favorable outcomes are $1,2$; $1,3$; $2,1$; and $3,1$, we have 4 favorable outcomes. We have 20 total outcomes. Probability will be  $\frac{1}{5}$. Hence the answer is $\boxed{D}$.

 



\textbf{Solution 2}

We only have to analyze the first two draws as that gives us insight into if a third draw is necessary. Also, note that it is necessary to draw a $1$ to have 3 draws, otherwise $5$ will be attainable in two or fewer draws. So the probability of getting a $1$ is $\frac{1}{5}$. It is necessary to pull either a $2$ or $3$ on the next draw and the probability of that is $\frac{1}{2}$. But, the order of the draws can be switched so we get:

 $\frac{1}{5} \cdot \frac{1}{2} \cdot 2 = \frac{1}{5}$, or $\boxed {D}$




\textbf{Solution 3 (Inequalities)}

We know that the first $2$ draws must be $\le{3}$ leaving us with with the only option being that $1$ and $2$ are chosen in either order. This means there is a probability of $\frac{1}{\binom52} \cdot 2 = \frac{1}{5}$, or $\boxed {D}$



	\item (\textit{AMC 10 10B 2018}) In the figure below, $N$ congruent semicircles are drawn along a diameter of a large semicircle, with their diameters covering the diameter of the large semicircle with no overlap. Let $A$ be the combined area of the small semicircles and $B$ be the area of the region inside the large semicircle but outside the small semicircles. The ratio $A:B$ is $1:18$. What is $N$?


 


                    \begin{figure}[H]
                        \centering
                        \includegraphics[width=0.5\textwidth]{Problem/AMC_10_10B_2018_7_Problem_Figure_1.png}
                    \end{figure}
                    
$\textbf{(A) }16 \qquad \textbf{(B) }17 \qquad \textbf{(C) }18 \qquad \textbf{(D) }19 \qquad \textbf{(E) }36 \qquad$

 \textbf{Answer: }D



\textbf{Solution 1 (Work using Answer Choices)}

Use the answer choices and calculate them. The one that works is $\bold{\boxed{\text{(D)19}}}$.

 



\textbf{Solution 2 (More Algebraic Approach)}

Let the number of semicircles be $n$ and let the radius of each semicircle be $r$. To find the total area of all of the small semicircles, we have $n \cdot \frac{\pi \cdot r^2}{2}$. 

 Next, we have to find the area of the larger semicircle. The radius of the large semicircle can be deduced to be $n \cdot r$. So, the area of the larger semicircle is $\frac{\pi \cdot n^2 \cdot r^2}{2}$. 

 Now that we have found the area of both A and B, we can find the ratio. $\frac{A}{B}=\frac{1}{18}$, so part-to-whole ratio is $1:19$. When we divide the area of the small semicircles combined by the area of the larger semicircles, we get $\frac{1}{n}$. This is equal to $\frac{1}{19}$. By setting them equal, we find that $n = 19$. This is our answer, which corresponds to choice $\bold{\boxed{\text{(D)19}}}$.

 Solution by: Archimedes15

 
Note: using WLOG you could have just set $r$ as $1$ in the beginning.

 



\textbf{Solution 3 (Fakesolve)}

Each small semicircle is $\frac{1}{N^2}$ of the large semicircle. Since $N$ small semicircles make $\frac{1}{19}$ of the large one, $\frac{N}{N^2} = \frac1{19}$. Solving this, we get $\boxed{19}$. 

 ... Don't think this one works but it kinda does

 


	\item (\textit{AMC 10 10B 2018}) Sara makes a staircase out of toothpicks as shown:


 
                    \begin{figure}[H]
                        \centering
                        \includegraphics[width=0.5\textwidth]{Problem/AMC_10_10B_2018_8_Problem_Figure_1.png}
                    \end{figure}
                    


 This is a $3$-step staircase and uses $18$ toothpicks. How many steps would be in a staircase that used $180$ toothpicks?


 $\textbf{(A) }10 \qquad \textbf{(B) }11 \qquad \textbf{(C) }12 \qquad \textbf{(D) }24 \qquad \textbf{(E) }30 \qquad$

 \textbf{Answer: }C



\textbf{Solution 1}

Notice that the sequence of numbers of toothpicks from the 1st step to the 3rd step is $4, 10, 18$

 We find that the first difference between 4 and 10 is 6 and 10 and 18 is 8.

 So, the second difference is constant, 2 (the difference between 6 and 8)

 Because the second differences are constant, the sequence comes from a quadratic polynomial in the form $ax^2+bx+c$

 Creating a system of equations, we find that (by plugging in values to $ax^2+bx+c$):

 $a + b + c = 4$$4a + 2b + c = 10$$9a + 3b + c = 18$

 So, (a,b,c) is (1,3,0) and we get the quadratic $x^2 + 3x$

 Set $x^2 + 3x = 180$ and factor

 $(x-12)(x+15)$

 So x = 12, and from this we see that the solution is $\boxed {(C) 12}$
\begin{verbatim}
By: UnknownAnsh
\end{verbatim}





\textbf{Solution 2}

A staircase with $n$ steps contains $4 + 6 + 8 + ... + 2n + 2$ toothpicks. This can be rewritten as $(n+1)(n+2) -2$. 

 So, $(n+1)(n+2) - 2 = 180$

 So, $(n+1)(n+2) = 182.$ 

 Inspection could tell us that $13 \cdot 14 = 182$, so the answer is $13 - 1 = \boxed {(C) 12}$




\textbf{Solution 3}

Layer $1$: $4$ steps

 Layer $1,2$: $10$ steps

 Layer $1,2,3$: $18$ steps

 Layer $1,2,3,4$: $28$ steps

 From inspection, we can see that with each increase in layer the difference in toothpicks between the current layer and the previous increases by $2$. Using this pattern:

 $4, 10, 18, 28, 40, 54, 70, 88, 108, 130, 154, 180$

 From this we see that the solution is $\boxed {(C) 12}$

 By: Soccer\_JAMS

 



\textbf{Solution 4 (10th Grade Math)}

Notice that we have the points \( (1,4) \), \( (2,10) \), \( (3,18) \). Then, we calculate the first difference to get 6 and 8 respectively, and calculate the second difference to get 2. This is a quadratic expression in the form \( ax^2 + bx + c \).

 Because we have a table of elements, \( 2a = \) the second difference, so \( 2a = 2 \Rightarrow a = 1 \). The quadratic is now \( y = x^2 + bx + c \). We plug in two points to get \( 4 = b + c \) and \( 10 = 4 + 2b + c \Rightarrow 6 = 2b + c \). Solving through elimination gives us \( b = 2 \), and \( c = 2 \).

 We now have \( y = x^2 + 2x + 2 \). We put \( y = 180 \), and see that \( 180 = x^2 + 2x + 2 \), and \( 0 = x^2 + 2x - 178 \). We use the quadratic formula to get \( -2 \pm \sqrt{716}/2 \). \( \sqrt{716} \) is greater than \( 26^2 \) but less that \( 27^2 \), so we say \( \sqrt{716} = 26 \). We then get 12 or -14.

 We cannot have a negative number of steps, so our answer is $\boxed {(C) 12}$.

 



\textbf{Solution 5}

We can find a function that gives us the number of toothpicks for every layer. Using finite difference, we know that the degree must be $2$ and the leading coefficient is $1$. The function is $f(n)=n^2+3n$ where $n$ is the layer and $f(n)$ is the number of toothpicks.

 
We have to solve for $n$ when $n^2+3n=180\Rightarrow n^2+3n-180=0$. Factor to get $(n-12)(n+15)$. The roots are $12$ and $-15$. Clearly $-15$ is impossible so the answer is $\boxed {(C) 12}$.

 



\textbf{Solution 6}

Notice that the number of toothpicks can be found by adding all the horizontal and all the vertical toothpicks. We can see that for the case of 3 steps, there are $2(3+3+2+1)=18$ toothpicks. Thus, the equation is $2S + 2(1+2+3...+S)=180$ with $S$ being the number of steps. Solving, we get $S = 12$, or $\boxed {(C) 12}$.-liu4505

 



\textbf{Solution 6 General Formula}

There are $\frac{n(n+1)}{2}$ squares. Each has $4$ toothpick sides. To remove overlap, note that there are $4n$ perimeter toothpicks. $\frac{\frac{n(n+1)}{2}\cdot 4-4n}{2}$ is the number of overlapped toothpicksAdd $4n$ to get the perimeter (non-overlapping). Formula is $\text{number of toothpicks} = n^2+3n$ Then you can "guess" or factor (also guessing) to get the answer $\boxed{\text{(C) }12}$.~bjc

 \textbf{Not a solution! Just an observation.} If you are trying to look for a pattern, you can see that the first column is made of 4 toothpicks. The second one is made from 2 squares: 3 toothpicks for the first square and 4 for the second. The third one is made up of 3 squares: 3 toothpicks for the first and second one, and 4 for the third one. The pattern continues like that. So for the first one, you have 0 "3 toothpick squares" and 1 "4 toothpick squares". The second is 1 to 1. The third is 2:1. And the amount of three toothpick squares increase by one every column.

 The list is as follow for the number of toothpicks used...$4$,$4+3$,$4+6$,$4+9$, and so on.4, 7, 10, 13, 16, 19, ...

 - Flutterfly

 


	\item (\textit{AMC 10 10B 2018}) The faces of each of $7$ standard dice are labeled with the integers from $1$ to $6$. Let $p$ be the probability that when all $7$ dice are rolled, the sum of the numbers on the top faces is $10$. What other sum occurs with the same probability $p$?


 $\textbf{(A) }13 \qquad \textbf{(B) }26 \qquad \textbf{(C) }32 \qquad \textbf{(D) }39 \qquad \textbf{(E) }42 \qquad$

 \textbf{Answer: }D



\textbf{Solution 1 (Casework; Using Answer Choices)}

The number 10 can be achieved by \( 4+6 \), \( 5+5 \), \( 6+4 \), in which there must exist at least 5 number ones. 

 Thus, we can write an answer in the form \( 1 + 1 + 1 + 1 + 1 + x + y \), where \( x + y \) must sum to 5. We notice that \( x + y = 5 \) can be either \( 1 + 4 \), \( 2 + 3 \), \( 3 + 2 \), or \( 4 + 1 \). So \( x + y = 4\) cases.

 The question now is which of the following answer choices can be written as a 5 term consecutive sum with the last two numbers \( x\) and \(y\) having exactly 4 cases.

 We first take E. 42. We see that we can have 5 consecutive number 6's to get 30. We then have a remaining 12, which can be represented as 6 + 6 only, not giving us our wanted 4 cases. 

 We try D. 39. We see we can write this as 5 consecutive number 6's. This gives us a remainder of 9, which can be written as 6+3, 4+5, 5+4, or 3+6, giving us 4 cases, exactly like the cases we got for 10.

 Our answer is therefore $\boxed{\textbf{(D)} \text{ 39}}$.

 



\textbf{Solution 1.1}

If you wish to check your answer.

 32 = a consecutive number 5's or 6's, gives a remainder of 7 and 2 respectively.

 The consecutive number of 6's obviously doesn't work.

 The consecutive number of 5's gives the remainder 7, which is \( 1+6 \), \( 2 + 5 \), \( 3+4 \), \( 4+3 \), \( 5+2 \), \( 6+1 \), which is 6 cases and not 4.

 26 is either consecutive 5's or 4's. 

 This gives a remainder of 1 or 6. 1 obviously doesn't work. 6 is \( 1+5 \), \( 2+4 \), \( 3+3 \), \( 4+2 \), or \( 5+1 \), which is 5 cases, not 4.

 13 is consecutive 2's or 1's, giving a remainder of 3 and 7 respectively. 

 We see that 3 doesn't work, and we already know that 7 doesn't work either.

 



\textbf{Solution 1.2}

A dice has the property all sums are symmetrically distributed (as the dice is fair). Therefore, one can realize that a sum of 10 (3 more than the minimum sum 7) comes the same as 3 less than the maximal sum 42, or $\boxed{39}$.

 Minor edits by Frogflick

 



\textbf{Solution 2}

It can be seen that the probability of rolling the smallest number possible is the same as the probability of rolling the largest number possible, the probability of rolling the second smallest number possible is the same as the probability of rolling the second largest number possible, and so on. This is because the number of ways to add a certain number of ones to an assortment of $7$ ones is the same as the number of ways to take away a certain number of ones from an assortment of $7$ $6$s.

 So, we can match up the values to find the sum with the same probability as $10$. We can start by noticing that $7$ is the smallest possible roll and $42$ is the largest possible roll. The pairs with the same probability are as follows:

 $(7, 42), (8, 41), (9, 40), (10, 39), (11, 38)...$

 However, we need to find the number that matches up with $10$. So, we can stop at $(10, 39)$ and deduce that the sum with equal probability as $10$ is $39$. So, the correct answer is $\boxed{\textbf{(D)} \text{39}}$, and we are done.

 Written By: Archimedes15

 


 Add-on by ike.chen: to see how the number of ways to roll $10$ and $39$ are the same, consider this argument:

 Each of the $7$ dice needs to have a nonnegative value; it follows that the number of ways to roll $10$ is $\binom {10-1}{7-1}=84$ by stars and bars. $10-7=3$, so there's no chance that any dice has a value $>6$.

 Now imagine $7$ piles with $6$ blocks each. The number of ways to take $3$ blocks away (making the sum $7\cdot 6-3=39$) is also $\binom {3+7-1}{7-1}=84$.

 correction to the add-on: the die need a positive value. We first give every die 1 so we have $10-7=3$ "balls" left to put into the 7 more bins/die. a dice can have a value of 0 for the \# of added balls. Thus, from stars and bars it follows that there are $\binom{7+3-1}{7-1}=84$ ways

 



\textbf{Solution 3}

Let's call the unknown value $x$. By symmetry, we realize that the difference between 10 and the minimum value of the rolls is equal to the difference between the maximum and $x$. So, 

 $10 - 7 = 42- x$
$x = 39$ and our answer is  $\boxed{\textbf{(D)} \text{ 39}}$By: Soccer\_JAMS

 



\textbf{Solution 4}

For the sums to have equal probability, the average sum of both sets of $7$ dies has to be $(6+1)\cdot 7 = 49$. Since having $10$ is similar to not having $10$, you just subtract 10 from the expected total sum. $49 - 10 = 39$ so the answer is $\boxed{\textbf{(D)} \text{ 39}}$

 By: epicmonster

 Revised solution by Williamgolly (includes bijections):Notice that we first must have at least a 1 on each die.  Now, we form the following bijection: biject each value on the original die, say $x$ to a value $7-x$ on the new die. Notice how now we need to 'take away' 3 from seven dies that all show 6.  Therefore, the answer is 39.

 Note: This is basically just having $10$ on the tops of the dice, and asking about the sum of the bottom of the dice. The opposite sides of a dice add up to $7$. ~Fib\_0112358

 



\textbf{Solution 5}

The expected value of the sums of the die rolls is $3.5\cdot7=24.5$, and since the probabilities should be distributed symmetrically on both sides of $24.5$, the answer is $24.5+(24.5-10)=39$, which is $\boxed{\textbf{(D)} \text{ 39}}$.

 By: dajeff

 



\textbf{Solution 6}

Another faster and easier way of doing this, without using almost any math at all, is realizing that the possible sums are ${7,8,9,10,...,39,40,41,42}$. By symmetry, (and doing a few similar problems in the past), you can realize that the probability of obtaining $7$ is the same as the probability of obtaining $42$, $P(8)=P(41)$ and on and on and on. This means that $P(10)=P(39)$, and thus the correct answer is $\boxed{\textbf{(D)} \text{ 39}}$.

 By: fhdsaukfaioifk

 \textbf{Note} Calculating the probability of getting a sum of $10$ is also easy. There are $3$ cases:

 
Case $1$: $\{1,1,1,1,1,1,4\}$
${7 \choose 6}=7$ cases

 
Case $2$: $\{1,1,1,1,1,2,3\}$
${7 \choose 5}=6 \cdot 7=42$ cases

 
Case $3$: $\{1,1,1,1,2,2,2\}$
${7 \choose 4}=\frac {7 \cdot 6 \cdot 5}{3 \cdot 2}=35$ cases

 
The probability is ${84 \over 6^7} = \frac{14}{6^6}$. 

 Calculating $6^6$: 

 $6^6=(6^3)^2=216^2=46656$ 

 Therefore, the probability is ${14 \over 46656} = \boxed{{7 \over 23328}}$

 \textbf{Related Problems} There is similar to problem 11 of the AMC 10A in the same year, which is almost a replica of the problem mentioned by Zeric Hang in the Note section:\href{https://artofproblemsolving.com/wiki/index.php/2018\_AMC\_10A\_Problems/Problem\_11}{https://artofproblemsolving.com/wiki/index.php/2018\_AMC\_10A\_Problems/Problem\_11}

 


	\item (\textit{AMC 10 10B 2018}) In the rectangular parallelepiped shown, $AB=3$, $BC=1$, and $CG=2$. Point $M$ is the midpoint of $\overline{FG}$. What is the volume of the rectangular pyramid with base $BCHE$ and apex $M$? 


 
                    \begin{figure}[H]
                        \centering
                        \includegraphics[width=0.5\textwidth]{Problem/AMC_10_10B_2018_10_Problem_Figure_1.png}
                    \end{figure}
                    
$\textbf{(A) }1 \qquad \textbf{(B) }\frac{4}{3} \qquad \textbf{(C) }\frac{3}{2} \qquad \textbf{(D) }\frac{5}{3} \qquad \textbf{(E) }2 \qquad$

 \textbf{Answer: }E



\textbf{Solution 1}

Consider the cross-sectional plane and label its area $b$. Note that the volume of the triangular prism that encloses the pyramid is $\frac{bh}{2}=3$, and we want the rectangular pyramid that shares the base and height with the triangular prism. The volume of the pyramid is $\frac{bh}{3}$, so the answer is $\boxed{(E) 2}$.(AOPS12142015)

 



\textbf{Solution 2}

We can start by finding the total volume of the parallelepiped. It is $2 \cdot 3 \cdot 1 = 6$, because a rectangular parallelepiped is a rectangular prism.

 Next, we can consider the wedge-shaped section made when the plane $BCHE$ cuts the figure. We can find the volume of the triangular pyramid with base $EFB$ and apex $M$. The area of $EFB$ is $\frac{1}{2} \cdot 2 \cdot 3 = 3$. Since $\overline{BC}$ is given to be $1$, we have that $\overline{FM}$ is $\frac{1}{2}$. Using the formula for the volume of a triangular pyramid, we have $V = \frac{1}{3} \cdot \frac{1}{2} \cdot 3 = \frac{1}{2}$. Also, since the triangular pyramid with base $HGC$ and apex $M$ has the exact same dimensions, it has volume $\frac{1}{2}$ as well.

 The original wedge we considered in the last step has volume $3$, because it is half of the volume of the parallelepiped. We can subtract out the parts we found to have $3 - \frac{1}{2} \cdot 2 = 2$. Thus, the volume of the figure we are trying to find is $2$. This means that the correct answer choice is $\boxed{E}$.

 Written by: Archimedes15

 NOTE: For those who think that it isn't a rectangular prism, please read the problem. It says "rectangular parallelepiped." If a parallelepiped is such that all of the faces are rectangles, it is a rectangular prism.

 



\textbf{Solution 3 (Competition Speed)}

Notice that the volume of the rectangular prism is \( lwh = 3 \cdot 1 \cdot 2 = 6 \)

 We first see that the triangular prism EHCBA is just half the volume of the rectangular prism which is just \( 6/2 = 3 \).

 We next notice that to get the pyramid with apex \( M \), we need to subtract another pyramid \( GHMC \) that is right behind the pyramid we wish to find. We notice the gap between the answer choices, and realize we can assume a formula for pyramid \( GHMC \).

 So we say that the formula is \( V = \frac{1}{3} Bh \), \( B \) is the base area and \( h \) is the height.

 The length \( MG \) is just \( 1/2 \), and to find \( B \) we do \( \frac{1}{2} \cdot \frac{1}{2} \cdot 2 = \frac{1}{2} \). The volume of pyramid \( GHMC \) is \( V = \frac{1}{3} \cdot \frac{1}{2} \cdot 3 = \frac{1}{2} \).

 We subtract the total volume 6 by the combined volume \( 3 + \frac{1}{2} = \frac{7}{2} \) to get \( \frac{12}{2} - \frac{7}{2} = \frac{5}{2} \). \( \frac{5}{2} \) is just 2.5, and the closest answer choice is $\boxed{\textbf{E } 2}$.

 



\textbf{Solution 4}

If you look carefully, you will see that on the either side of the pyramid in question, there are two congruent tetrahedra. The volume of one is $\frac{1}{3}Bh$, with its base being half of one of the rectangular prism's faces and its height being half of one of the edges, so its volume is $\frac{1}{3} (3 \times 2/2 \times \frac{1}{2})=\frac{1}{2}$. We can obtain the answer by subtracting twice this value from the diagonal half prism, or$(\frac{1}{2} \times 3 \times 2 \times 1) - (2 \times \frac{1}{2})=$ $\boxed{2}$




\textbf{Solution 5}

You can calculate the volume of the rectangular pyramid by using the formula, $\frac{Ah}{3}$. $A$ is the area of the base, $BCHE$, and is equal to $BC * BE$. The height, $h$, is equal to the height of triangle $FBE$ drawn from $F$ to $BE$.

 $BE=\sqrt{BF^2 + EF^2}=\sqrt{13}$ Area of $BCHE = BC * BE = \sqrt{13}$
$h = 2 *$ Area of $FBE / BE$ (since Area $= \frac{1}{2}bh$).

 Area of $FBE = \frac{1}{2} * FB * FE = 3$
$h = 2 * 3 / \sqrt{13} = \frac{6}{\sqrt{13}}$

 Volume of pyramid $=\frac{1}{3} * \sqrt{13} * \frac{6}{\sqrt{13}} = 2$

 Answer is $\boxed{\textbf{E } 2}$




\textbf{Solution 6}

We can start by identifying the information we need. We need to find the area of rectangle $EHCB$ and the height of rectangular prism $EHCBM$.

 In order to find the area of $EHCB,$ we can use the Pythagorean Theorem. We find that $EB = \sqrt{13}$, so the area of rectangle $EHCB = \sqrt{13}$. We shall refer to this as $x$.

 In order to find the height of rectangular prism $EHCBM$, we can examine triangle $EFB$. We can use the Geometric Mean Theorem to find that when an altitude is dropped from point $F,$ $\overline{EB}$ is split into segments of length $\dfrac{4 \cdot \sqrt{13}}{13}$ and $\dfrac{9 \cdot \sqrt{13}}{13}$. Taking the geometric mean of these numbers, we find that the altitude has length $\dfrac{6 \cdot \sqrt{13}}{13}$. This is also the height of the rectangular prism, which we shall refer to as $y$.

 Plugging $x$ and $y$ into the formula $V = \dfrac{b \cdot h}{3},$ we find that the volume is $\boxed{2}$. The answer is $\boxed{E}$.

 



\textbf{Solution 7}

We start by setting the formula for the volume of a rectangular pyramid: $\frac{1}{3}Bh$. By the Pythagorean Theorem, we know that $BE = \sqrt{BF^2 + EF^2} = \sqrt{13}$. Therefore, the area of the base is $1 \times \sqrt{13} = \sqrt{13}$. Next, we would like to know the height of the pyramid. We can observe that the altitude from point $F$ in $\triangle EFB$ is parallel to the height of the pyramid and therefore congruent because those two altitudes are on the same plane of base $EBCH$. From this, we only need to find the altitude from point $F$ in $\triangle EFB$ and plug it into our formula for the volume of a rectangular pyramid. This is easy because we already know the area of $\triangle EFB$ and the base from point $F$, so all we need to do is divide: $\frac{2 \times 3}{\sqrt{13}} = \frac{6}{\sqrt{13}} = \frac{6\sqrt{13}}{13}$. Now all we need to do is plug in all our known values into the volume formula: $\frac{1}{3}Bh = \frac{\sqrt{13} \times \frac{6\sqrt{13}}{13}}{3} = \boxed{(E) 2}$




\textbf{Solution 8}


                    \begin{figure}[H]
                        \centering
                        \includegraphics[width=0.5\textwidth]{Solution/AMC_10_10B_2018_10_Solution_Figure_1.jpg}
                    \end{figure}
                    

 Using the Pythagorean Theorem, we can easily find that $EB = \sqrt{2^2 + 3^2} = \sqrt{13}$. Quickly computing, we find the area of the base, $BCHE = \sqrt{13} \cdot 1 = \sqrt{13}$. Now we can make the following adjustments to our 3d shape as shown in the diagram. All we need now is to solve for the height, or $XM$. We can set up to following equation due to our knowledge of altitudes(of the hypotenuse)in right triangles. We can set up the following equations:
 \begin{align*} b(a+b) &= (MH_1)^2 \\ a(a+b) &= (MH_2)^2 \\ b\sqrt{13} &= 3^2 \\ a\sqrt{13} &= 2^2 \\ b &= \dfrac{9}{\sqrt{13}} \\ a &= \dfrac{4}{\sqrt{13}} \\ (MX)^2 &= ab \\ (MX)^2 &= \dfrac{9 \cdot 4}{13} \\ MX &= \dfrac {3 \cdot 2}{\sqrt{13}} \\ \end{align*}
Thus $\triangle V_{BCHEM} = \dfrac{\text{(height)}\cdot \text{(base)}}{3} = \dfrac{MX \cdot BCHE}{3} = \dfrac {\sqrt{13} \cdot \dfrac{3 \cdot 2}{\sqrt{13}}}{3}$ $= \boxed{\textbf{(E) } 2}$




\textbf{Solution 9 (Basic Coordinate Geometry)}

First of all, we can quickly calculate the base area $EBCH$ to be $\sqrt{13}$, since $BC = 1$ (given) and $EB = \sqrt{2^2+3^2} = \sqrt{13}$ (Pythagorean Theorem).

 Now all we need is to find the height of $EBCHM$, which we can denote as $h$. We know that the length of the altitude from $M$, $F$, and $G$ to base $EBCH$ is $h$, and that's the length that will be easiest to use to find the height. To make this problem simpler, you can turn this 3D model into a 2D one (look at face $EABF$).

 Looking at face $EABF$, $h$ would be equal to the length of the altitude from $F$ to $EB$. To find this, we can find the coordinates of $F$ and the point where the altitude from $F$ touches $EB$, which we will call point $X$. Immediately, we know that the coordinates of $F$ are $(3,2)$, given the values on the rectangle, and we can find the coordinates of point $X$ by first noticing that the equation of $EB$ is $y = \dfrac{-2x}{3}+2$. Because of perpendicular slopes being the inverse reciprocal of the original slope, the equation for FX is $y = \dfrac{3x}{2}-\dfrac{5}{2}$. From this, $y = \dfrac{3x}{2}+b$, and $2=\dfrac{3(3)}{2}+b$, so $b = \dfrac{-5}{2}$. Now the intersection of these two equations can be found to be $(\dfrac{27}{13},\dfrac{8}{13})$. We can use the distance formula to find the distance between the two points. $\sqrt{(3-\dfrac{27}{13})^2 + (2-\dfrac{8}{13})^2} = \sqrt{(\dfrac{12}{13})^2+(\dfrac{8}{13})^2} = \dfrac{\sqrt{468}}{13}$. 

 Now that we know the height length and base area of the pyramid $EBCHM$, we can use the volume formula for it: $\dfrac{\text{(base area)} \cdot \text{(height length)}}{3} = \dfrac{\dfrac{\sqrt{468}}{13} \cdot \sqrt{13}}{3}$ $= \boxed{\textbf{(E) } 2}$




\textbf{Solution 10 (Point-line distance formula)}

In the rectangular parallelepiped, we are given
 \[AB = 3,\quad BC = 1,\quad CG = 2.\]The base of the pyramid is the quadrilateral $BCHE$.

 Since opposite edges are perpendicular, $BCHE$ is a rectangle with
 \[BC = 1 \quad \text{and} \quad BE = \sqrt{AB^2 + AE^2} = \sqrt{3^2 + 2^2} = \sqrt{13}.\]So the area is
 \[[BCHE] = BC \cdot BE = 1 \cdot \sqrt{13} = \sqrt{13}.\]

 Project the solid into the 2D plane containing $B, E,$ and $M$. Use the coordinates
 \[B = (3,0),\quad E = (0,2),\quad M = (3,2).\]The slope of $\overline{BE}$ is
 \[\frac{2-0}{0-3} = -\frac{2}{3},\]so its equation is
 \[y = -\frac{2}{3}x + 2.\]Rewriting in standard form:
 \[2x + 3y - 6 = 0.\]

 Using the point-to-line distance formula, the perpendicular distance from $M = (3,2)$ to this line is
 \[\frac{|2(3) + 3(2) - 6|}{\sqrt{2^2 + 3^2}} = \frac{6}{\sqrt{13}}.\]This is the height of the pyramid.

 The volume of a pyramid is
 \[V = \frac{1}{3} \cdot (\text{base area}) \cdot (\text{height}).\]Thus,
 \[V = \frac{1}{3} \cdot \sqrt{13} \cdot \frac{6}{\sqrt{13}} = \frac{6}{3} = 2.\]



	\item (\textit{AMC 10 10B 2018}) Which of the following expressions is never a prime number when $p$ is a prime number?


 $\textbf{(A) }p^2+16 \qquad \textbf{(B) }p^2+24 \qquad \textbf{(C) }p^2+26 \qquad \textbf{(D) }p^2+46 \qquad \textbf{(E) }p^2+96 \qquad$

 \textbf{Answer: }C



\textbf{Solution 1}

Each expression is in the form $p^2 + n$.

 All prime numbers are of the form $6k \pm 1$, AKA they are congruent to $\pm1 \pmod{6}$. We can utilize this nicely to check for what we are looking for. If the expression is a prime, then

 
$p^2 + n \equiv \pm1 \pmod{6}$
$\Rightarrow (\pm1)^2 + n \equiv \pm1 \pmod{6}$
$\Rightarrow 1 + n \equiv \pm1 \pmod{6}$
$\Rightarrow n \equiv (\pm1) - 1 \pmod{6}$
$\Rightarrow n \equiv (1$ $or$ $-1) - 1 \pmod{6}$
$\Rightarrow n \equiv 0$ $or$ $-2 \pmod{6}$

 
Now, just check for $n$ in each option using this condition to check whether its prime or not.

 
 \[(A)\ n = 16; prime\]
 \[(B)\ n = 24; prime\]
 \[(C)\ n = 26; not\ prime\]
 \[(D)\ n = 46; prime\]
 \[(E)\ n = 96; prime\]

 Therefore, the answer is $\boxed{\textbf{(C) } p^2 + 26}$.

 


 



\textbf{Solution 2}

Because squares of a non-multiple of 3 is always $1 \pmod{3}$, the only expression always a multiple of $3$ is $\boxed{\textbf{(C) } p^2+26}$. This is excluding when $p\equiv 0 \pmod{3}$, which only occurs when $p=3$, then $p^2+26=35$ which is still composite.

 



\textbf{Solution 3}

\( p^2 + 16 \) is prime for \( p = 5 \), so we know that \( A \) doesn't work.

 Because \( p^2 + 96 = p^2 + (16) \cdot 6 \), we see that eventually there exists a number \( p \) that makes the sum prime, therefore \( E \) also doesn't work.

 Option \( B \) is \( p^2 + 24 \), but because the number itself has no carryover, we can assume \( p = 5 \) again to get the number \( 41 \), which is prime.

 Option \( D \) is \( p^2 + 46 \), ending in \( 6 \), so we try \( p = 5 \) again to get \( 71 \). This is also prime.

 We can either deduce through elimination that \( p^2 + 26 \) is always composite for any prime \( p \), but we can prove it by showing \( p = 5 \) doesn't work, \( p = 7 \) doesn't work, and \( p = 11 \) doesn't work, and through \href{https://artofproblemsolving.com/wiki/index.php/Induction?srsltid=AfmBOooY-SKw1NAbaPWjFKkyfSoHj6RXiwWuqty6FtGzCI5GAw7eyCTA}{Engineer's Induction}, \( p^2 + 26 \) always remains composite.

 Therefore the answer is $\boxed{\textbf{(C) } p^2 + 26}$.

 



\textbf{Solution 4 (Answer Choices)}

Since the question asks which of the following will never be a prime number when $p$ is a prime number, a way to find the answer is by trying to find a value for $p$ such that the statement above won't be true.

 A) $p^2+16$ isn't true when $p=5$ because $25+16=41$, which is prime

 B) $p^2+24$ isn't true when $p=7$ because $49+24=73$, which is prime

 C) $p^2+26$

 D) $p^2+46$ isn't true when $p=5$ because $25+46=71$, which is prime

 E) $p^2+96$ isn't true when $p=19$ because $361+96=457$, which is prime

 Therefore, $\framebox{C}$ is the correct answer.

 -DAWAE

 


 Minor edit by Lucky1256. P=\_\_\_ was the wrong number.

 More minor edits by beanlol.

 More minor edits by mathmonkey12.

 


	\item (\textit{AMC 10 10B 2018}) Line segment $\overline{AB}$ is a diameter of a circle with $AB=24$. Point $C$, not equal to $A$ or $B$, lies on the circle. As point $C$ moves around the circle, the centroid (center of mass) of $\triangle{ABC}$ traces out a closed curve missing two points. To the nearest positive integer, what is the area of the region bounded by this curve?


 $\textbf{(A) }25 \qquad \textbf{(B) }38 \qquad \textbf{(C) }50 \qquad \textbf{(D) }63 \qquad \textbf{(E) }75 \qquad$

 \textbf{Answer: }C



\textbf{Solution 1}

By the Inscribed Angle Theorem, $\triangle ABC$ is a right triangle with $\angle C=90^{\circ}.$ So, its circumcenter is the midpoint of $\overline{AB},$ and its median from $C$ is half as long as $\overline{AB}.$ For each $\triangle ABC,$ let $O$ and $G$ be its circumcenter and centroid, respectively. It follows that $OA=OB=OC=12.$ In any triangle, since the centroid divides each median into parts in the ratio $2:1,$ with the centroid being twice as close to the midpoint of a side as it is to the opposite vertex, we have $OG=\frac13 OC=4.$

 As shown below, $\triangle ABC_1$ and $\triangle ABC_2$ are two shapes of $\triangle ABC$ with centroids $G_1$ and $G_2,$ respectively:
                    \begin{figure}[H]
                        \centering
                        \includegraphics[width=0.5\textwidth]{Solution/AMC_10_10B_2018_12_Solution_Figure_1.png}
                    \end{figure}
                    Therefore, point $G$ traces out a circle (missing two points) with the center $O$ and the radius $\overline{OG},$ as indicated in red. To the nearest positive integer, the area of the region bounded by the red curve is $\pi\cdot OG^2=16\pi\approx\boxed{\textbf{(C) } 50}.$




\textbf{Solution 2}

We assign coordinates. Let $A = (-12,0)$, $B = (12,0)$, and $C = (x,y)$ lie on the circle $x^2 +y^2 = 12^2$. Then, the centroid of $\triangle ABC$ is $G = \left(\frac{-12 + 12 + x}{3}, \frac{0 + 0 + y}{3}\right) = \left(\frac x3,\frac y3\right)$. Thus, $G$ traces out a circle with a radius $\frac13$ of the radius of the circle that point $C$ travels on. Thus, $G$ traces out a circle of radius $\frac{12}{3} = 4$, which has area $16\pi\approx \boxed{\textbf{(C) } 50}$.

 



\textbf{Solution 3}

First we can draw a few conclusions from the given information. Firstly we can see clearly that the distance from the centroid to the center of the circle will remain the same no matter $C$ is on the circle. Also we can see that because the two legs of every triangles will always originate on the diameter, using inscribed angle rules, we know that $\angle C = \frac{180^\circ}{2} = 90^\circ$. Now we know that all triangles $ABC$ will be a right triangle. We also know that the closed curve will simply be a circle with radius equal to the centroid of each triangle. We can now pick any arbitrary triangle, calculate its centroid, and the plug it in to the area formula. Using a $45^\circ$-$45^\circ$-$90^\circ$ triangle in conjunction with the properties of a centroid, we can quickly see that the length of the centroid is $4$ now we can plug it in to the area formula where we get $16\pi\approx\boxed{\textbf{(C) } 50}$.

 


 



\textbf{Solution 4}

Let the center of the circle be $O$. Since $AB$ is a diameter of length $24$, $OC$ is a median to $AB$ and a radius of the circle, so $OC = 12$. 

 Let the centroid of $\triangle ABC$ be $G$. By the Centroid Theorem, $G$ lies on the segment $OC$ such that the ratio of $OG$ to $GC$ is $1:2$. 

 Since $OG = \frac{1}{3} OC$, we have:
 \[OG = \frac{1}{3}(12) = 4\]

 As $C$ moves along the boundary of the circle, the centroid $G$ carves out a smaller, concentric circle of radius $4$. 

 The area of this region is:
 \[\text{Area} = \pi \cdot 4^2 = 16\pi \approx 50.27\]

 Rounding to the nearest whole number, the answer is $\boxed{50}$.

 



\textbf{Solution 5}

By barycentric coordinates, $G$ the centroid satisfies $G = \frac{a+b+c}{3}$ where $a,b,c$ are the position vectors of $A,B,C$ respectively. Then, by definition, let $a \langle 0,0 \rangle$ and $b \langle 24, 0 \rangle$. Then, $G = \langle 8,0 \rangle + \frac{c}{3}$. As $C$ lies on the circle, the vector value of $G$ is solely dependent on $\frac{c}{3}$. That is, as $C$ can lie on the circle, $\frac{c}{3}$ oscillates between $\langle 8,0 \rangle$ and $\langle 16,0 \rangle$, on the circle. Anywhere $C$ traverses, that value of $G$ is correspondent, hence if $C$ moves along the circle, then $G$ does as well along that interval. The requested path is then a circle concentric to circle $ABC$ whose radius is 4. That area is then $16 \cdot \pi \approx 16 \cdot 3 = 48$, so the nearest answer choice is $\boxed{50}$.

 


	\item (\textit{AMC 10 10B 2018}) How many of the first $2018$ numbers in the sequence $101, 1001, 10001, 100001, \dots$ are divisible by $101$?


 $\textbf{(A) }253 \qquad \textbf{(B) }504 \qquad \textbf{(C) }505 \qquad \textbf{(D) }506 \qquad \textbf{(E) }1009 \qquad$

 \textbf{Answer: }C



\textbf{Solution 1}

The number $10^n+1$ is divisible by 101 if and only if $10^n\equiv -1\pmod{101}$. We note that $(10,10^2,10^3,10^4)\equiv (10,-1,-10,1)\pmod{101}$, so the powers of 10 are 4-periodic mod 101.

 It follows that $10^n\equiv -1\pmod{101}$ if and only if $n\equiv 2\pmod 4$.

 In the given list, $10^2+1,10^3+1,10^4+1,\dots,10^{2019}+1$, the desired exponents are $2,6,10,\dots,2018$, and there are $\dfrac{2020}{4}=\boxed{\textbf{(C) } 505}$ numbers in that list.


 



\textbf{Solution 2 (Similar to solution 1)}

Note that $10^{2k}+1$ for some odd $k$ will satisfy $10^{2k}+1 \equiv -1\pmod {101}$. Each $2k \in \{2,6,10,\dots,2018\}$, so the answer is $\boxed{\textbf{(C) } 505}$




\textbf{Solution 3}

If we divide each number by $101$, we see a pattern occuring in every 4 numbers. $101, 1000001, 10000000001, \dots$. We divide $2018$ by $4$ to get $504$ with $2$ left over. Looking at our pattern of four numbers from above, the first number is divisible by $101$. This means that the first of the $2$ left over will be divisible by $101$, so our answer is $\boxed{\textbf{(C) } 505}$.

 



\textbf{Solution 4}

Note that $909$ is divisible by $101$, and thus $9999$ is too. We know that $101$ is divisible and $1001$ isn't so let us start from $10001$. We subtract $9999$ to get 2. Likewise from $100001$ we subtract, but we instead subtract $9999$ times $10$ or $99990$ to get $11$. We do it again and multiply the 9's by $10$ to get $101$. Following the same knowledge, we can use mod $101$ to finish the problem. The sequence will just be subtracting 1, multiplying by 10, then adding 1. Thus the sequence is ${0,-9,-99 ( 2),11, 0, ...}$. Thus it repeats every four. Consider the sequence after the 1st term and we have 2017 numbers. Divide $2017$ by four to get $504$ remainder $1$. Thus the answer is $504$ plus the 1st term or $\boxed{\textbf{(C) } 505}$.

 



\textbf{Solution 5}

Note that $101=x^2+1$ and $100...0001=x^n+1$, where $x=10$. We have that $\frac{x^n+1}{x^2+1}$ must have a remainder of $0$. By the remainder theorem, the roots of $x^2+1$ must also be roots of $x^n+1$. Plugging in $i,-i$ to $x^n+1$ yields that $n\equiv2\mod{4}$. Because the sequence starts with $10^2+1$, the answer is $\lceil 2018/4 \rceil=\boxed{\textbf{(C) } 505}$




\textbf{Solution 6 (Complex Numbers)}

The $n$th term of the sequence is $10^n+1$ for $n\ge 2$ (since $10^2+1=101$, $10^3+1=1001$, etc.). We seek those $n$ for which 
 \[101 \mid (10^n+1).\]Let $x = 10$. Note that $101=10^2+1$, so this is equivalent to the polynomial divisibility

 
 \[x^2+1 \mid x^n+1\]

 Over $\mathbb{C}$, $x^2+1$ has roots $i,-i$. Thus $x^2+1 \mid x^n+1$ if both $i,-i$ are roots of $x^n+1$, since,$i^n=-1 \quad \text{and} \quad (-i)^n=-1.$Because $i^4=1$, we have $i^n=-1$ exactly when $n\equiv 2 \pmod{4}$. For such $n$, $(-i)^n=(-1)^n i^n=i^n=-1$ as well, so the condition holds precisely when$n \equiv 2 \pmod{4}.$

 Among the first $2018$ terms of the sequence (which correspond to $n=2,3,\dots,2019$), the values of $n$ congruent to $2 \pmod{4}$ are$2,6,10,\dots,2018$.This is an arithmetic progression with common difference $4$; the count is

 $\frac{2018-2}{4}+1=\frac{2016}{4}+1=504+1=505$.Therefore, exactly $\boxed{505}$ of the first $2018$ numbers are divisible by $101$, which corresponds to choice $\boxed{\textbf{(C)}}$.

 


	\item (\textit{AMC 10 10B 2018}) A list of $2018$ positive integers has a unique mode, which occurs exactly $10$ times. What is the least number of distinct values that can occur in the list?


 $\textbf{(A) }202 \qquad \textbf{(B) }223 \qquad \textbf{(C) }224 \qquad \textbf{(D) }225 \qquad \textbf{(E) }234 \qquad$

 \textbf{Answer: }D



\textbf{Solution 1 (Statistics)}

To minimize the number of distinct values, we want to maximize the number of times a number appears. So, we could have $223$ numbers appear $9$ times, $1$ number appear once, and the mode appear $10$ times, giving us a total of $223 + 1 + 1 = \boxed{\textbf{(D)}\ 225}.$




\textbf{Solution 2 (Algebra)}

As in Solution 1, we want to maximize the number of time each number appears to do so. We can set up an equation $10 + 9( x - 1 )\geq2018,$ where $x$ is the number of values. Notice how we can then rearrange the equation into $1 + 9 ( 1 )+9 ( x - 1 )\geq2018,$ which becomes $9 x\geq2017,$ or $x\geq224\frac19.$ We cannot have a fraction of a value so we must round up to $\boxed{\textbf{(D)}\ 225}.$




\textbf{Solution 3}

There is a unique number \( n \) with a mode of 10. This means we have 2008 terms left. We see that \( \lfloor 2008/9 \rfloor = 223 \), with remainder 1. Our answer is simply \( 223 + 1 + 1 =\) $\boxed{\textbf{(D)}\ 225}.$



	\item (\textit{AMC 10 10B 2018}) A closed box with a square base is to be wrapped with a square sheet of wrapping paper. The box is centered on the wrapping paper with the vertices of the base lying on the midlines of the square sheet of paper, as shown in the figure on the left. The four corners of the wrapping paper are to be folded up over the sides and brought together to meet at the center of the top of the box, point $A$ in the figure on the right. The box has base length $w$ and height $h$. What is the area of the sheet of wrapping paper?


 
                    \begin{figure}[H]
                        \centering
                        \includegraphics[width=0.5\textwidth]{Problem/AMC_10_10B_2018_15_Problem_Figure_1.png}
                    \end{figure}
                    
$\textbf{(A) } 2(w+h)^2 \qquad \textbf{(B) } \frac{(w+h)^2}2 \qquad \textbf{(C) } 2w^2+4wh \qquad \textbf{(D) } 2w^2 \qquad \textbf{(E) } w^2h$

 \textbf{Answer: }A



\textbf{Solution 0, super simple and quick}

You could see that half the diagonal of the square wrapping paper is 
 \[\frac{w}{2} + h + \frac{w}{2} = w + h\]. Thus the full diagonal of the square wrapping paper would be 
 \[d = 2(w + h)\]. So the area would be 
 \[A = \frac{d^2}{2} = \frac{[2(w + h)]^2}{2} = \frac{4(w + h)^2}{2} = 2(w + h)^2\]. Thus the answer is 
 \[\boxed{\text{(A) } 2(w+h)^2}\]




\textbf{Solution 1}

Consider one-quarter of the image (the wrapping paper is divided up into 

 


	\item (\textit{AMC 10 10B 2018}) Let $a_1,a_2,\dots,a_{2018}$ be a strictly increasing sequence of positive integers such that 
 \[a_1+a_2+\cdots+a_{2018}=2018^{2018}.\] What is the remainder when $a_1^3+a_2^3+\cdots+a_{2018}^3$ is divided by $6$?


 $\textbf{(A) }0 \qquad \textbf{(B) }1 \qquad \textbf{(C) }2 \qquad \textbf{(D) }3 \qquad \textbf{(E) }4 \qquad$

 \textbf{Answer: }E



\textbf{Solution 1}

Verify that $a^3 \equiv a \pmod{6}$ manually for all $a\in \mathbb{Z}/6\mathbb{Z}$. We check: $0^3 \equiv 0 \pmod{6}$, $1^3 \equiv 1 \pmod{6}$, $2^3 \equiv 8 \equiv 2 \pmod{6}$, $3^3 \equiv 27 \equiv 3 \pmod{6}$, $4^3 \equiv 64 \equiv 4 \pmod{6}$, and $5^3 \equiv 125 \equiv 5 \pmod{6}$. We conclude that $a^3 \equiv a \pmod{6}$.

 Therefore, 
 \[a_1^3+a_2^3+\cdots+a_{2018}^3 \equiv a_1+a_2+\cdots+a_{2018}\pmod{6}.\]

 Thus the answer is congruent to $2018^{2018} = 1009^{2018} \cdot 2^{2018} \equiv \boxed{ \text{(E)}4} \pmod{6}$ because $2^n \pmod{6}$ alternates with $2$ and $4$ when $n$ increases, and by \href{https://artofproblemsolving.com/wiki/index.php?title=Euler\%27s\_Totient\_Theorem}{Euler's Totient Theorem} $1009^{2018} = (1009^2)^{1009} \equiv 1^{1009} \equiv 1 \pmod{6}$.

 



\textbf{Solution 2}

Note that $\left(a_1+a_2+\cdots+a_{2018}\right)^3=a_1^3+a_2^3+\cdots+a_{2018}^3+3a_1^2\left(a_1+a_2\\ +\cdots+a_{2018}-a_1\right)+3a_2^2\left(a_1+a_2+\cdots+a_{2018}-a_2\right)+\cdots+3a_{2018}^2\left(a_1+a_2+\cdots+a_{2018}-a_{2018}\right)+6\sum_{i\neq j\neq k}^{2018} a_ia_ja_k$

 Note that $a_1^3+a_2^3+\cdots+a_{2018}^3+3a_1^2\left(a_1+a_2+\cdots+a_{2018}-a_1\right)+3a_2^2\left(a_1+a_2+\cdots+a_{2018}-a_2\right)+\cdots+3a_{2018}^2\left(a_1+a_2+\cdots+a_{2018}-a_{2018}\right)+6\sum_{i\neq j\neq k}^{2018} a_ia_ja_k\equiv a_1^3+a_2^3+\cdots+a_{2018}^3+3a_1^2({2018}^{2018}-a_1)+3a_2^2({2018}^{2018}-a_2)+\cdots+3a_{2018}^2({2018}^{2018}-a_{2018}) \equiv -2(a_1^3+a_2^3+\cdots+a_{2018}^3)\pmod 6$Therefore, $-2(a_1^3+a_2^3+\cdots+a_{2018}^3)\equiv \left(2018^{2018}\right)^3\equiv\left( 2^{2018}\right)^3\equiv 4^3\equiv 4\pmod{6}$.

 Thus, $a_1^3+a_2^3+\cdots+a_{2018}^3\equiv 1\pmod 3$. However, since cubing preserves parity, and the sum of the individual terms is even, the sum of the cubes is also even, and our answer is $\boxed{\text{(E) }4}$




\textbf{Solution 3 (Partial Proof)}

First, we can not assume that the problem will have a consistent answer for all possible values of $a_1$. For the purpose of this solution, we will assume that $a_1 = 1$.

 We first note that $1^3+2^3+...+n^3 = (1+2+...+n)^2$. So what we are trying to find is what $\left(2018^{2018}\right)^2=\left(2018^{4036}\right)$ mod $6$. We start by noting that $2018$ is congruent to $2 \pmod{6}$. So we are trying to find $\left(2^{4036}\right) \pmod{6}$. Instead of trying to do this with some number theory skills, we could just look for a pattern. We start with small powers of $2$ and see that $2^1$ is $2$ mod $6$, $2^2$ is $4$ mod $6$, $2^3$ is $2$ mod $6$, $2^4$ is $4$ mod $6$, and so on... So we see that since $\left(2^{4036}\right)$ has an even power, it must be congruent to $4 \pmod{6}$, thus giving our answer $\boxed{\text{(E) }4}$. You can prove this pattern using mods. But I thought this was easier.

 -TheMagician

 



\textbf{Solution 4 (Lazy solution)}

First, we can assume that the problem will have a consistent answer for all possible values of $a_1$. For the purpose of this solution, assume $a_1, a_2, ... a_{2017}$ are multiples of 6 and find $2018^{2018} \pmod{6}$ (which happens to be $4$). Then ${a_1}^3 + ... + {a_{2018}}^3$ is congruent to $64 \pmod{6}$ or just $\boxed{\textbf{(E)}  4}$. 

 -Patrick4President

 WARNING: Technique does not always work; you have been warned

 



\textbf{Solution 5 (Even Lazier Solution)}

Due to the large amounts of variables in the problem, and the fact that the test is only 75 minutes, you can assume that the answer is probably just $2018^{2018} \pmod{6}$, which is $\boxed{\textbf{(E)}  4}$.

 \textbf{Algebraic Insight into Given Property} Mods is a good way to prove $a^3 \equiv a \pmod6$: residues are simply $3, \pm 2, \pm 1, 0$. Only $2$ and $3$ are necessary to check.Another way is to observe that $a^3-a$ factors into $(a-1)a(a+1)$. Any $k$ consecutive numbers must be a multiple of $k$, so $a^3-a$ is both divisible by $2$ and $3$. This provides an algebraic method for proving $a^3 \equiv a \pmod6$ for all $a$.

 



\textbf{Solution 6}

We are not given any other requirements than the two given in the problem, so we can assume that $a_1 + a_{2018} = a_2 + a_{2017} + \cdots + a_{1004} + a_{1005} = \frac{2018^{2018}}{2018} \cdot 2$. Since the (theoretical) $a_{1004.5}$ averages out to be $2018^{2017}$ (lets call this constant $N$), we can let $(a_1, a_{2018}) = (N - x_1, N+x_1)$, $(a_2, a_{2017}) = (N - x_2, N+x_2)$, ... ,$(a_{1004}, a_{1005}) = (N - x_{1004}, N+x_{1004})$ for the sequence of decreasing positive integers $x_1, x_2, \cdots, x_{1004}$. Notice that for any such pair (say $(a_{n}, a_{2019-n}$), we have $(a_n)^3 + (a_{2019-n})^3 = (N - x_n)^3 + (N+x_n)^3 = 2N^3 + 6N\cdot x_n \equiv 2N^3 \pmod{6}$. Thus, $a_1^3 + a_2^3 + \cdots + a_{2018}^3 \equiv 1004\cdot 2N^3 \equiv 2018(2018^{2017})^3 \equiv 2018^{6052} \equiv 2^{6052} \pmod{6}$. Since $2^x \pmod{6}$ repeats in the pattern $2, 4, 2, 4, \cdots$, our final answer is $\boxed{\textbf{(E)}  4}$




\textbf{Solution 7}

First, note that n\textasciicircum{}3 \equiv n (mod 6) for any integer n. This is because n\textasciicircum{}3 - n = (n-1)n(n+1) is the product of three consecutive integers, ensuring divisibility by both 2 and 3.

 Thus, Σ a\_i\textasciicircum{}3 \equiv Σ a\_i = 2018\textasciicircum{}2018 (mod 6).

 Since 2018 \equiv 2 (mod 6), we have 2018\textasciicircum{}2018 \equiv 2\textasciicircum{}2018 (mod 6).Observing the powers of 2 modulo 6:2\textasciicircum{}1 \equiv 2, 2\textasciicircum{}2 \equiv 4, 2\textasciicircum{}3 \equiv 2, 2\textasciicircum{}4 \equiv 4.The pattern alternates: even exponents yield 4. Since 2018 is even, 2\textasciicircum{}2018 \equiv 4 (mod 6).

 $\boxed{\textbf{(E)} \ 4}$



	\item (\textit{AMC 10 10B 2018}) In rectangle $PQRS$, $PQ=8$ and $QR=6$. Points $A$ and $B$ lie on $\overline{PQ}$, points $C$ and $D$ lie on $\overline{QR}$, points $E$ and $F$ lie on $\overline{RS}$, and points $G$ and $H$ lie on $\overline{SP}$ so that $AP=BQ<4$ and the convex octagon $ABCDEFGH$ is equilateral. The length of a side of this octagon can be expressed in the form $k+m\sqrt{n}$, where $k$, $m$, and $n$ are integers and $n$ is not divisible by the square of any prime. What is $k+m+n$?


 $\textbf{(A) }1 \qquad \textbf{(B) }7 \qquad \textbf{(C) }21 \qquad \textbf{(D) }92 \qquad \textbf{(E) }106 \qquad$

 \textbf{Answer: }B



\textbf{Solution 1}

Let $AP=BQ=x$. Then $AB=8-2x$.

 Now notice that since $CD=8-2x$ we have $QC=DR=x-1$.

 Thus by the Pythagorean Theorem we have $x^2+(x-1)^2=(8-2x)^2$ which becomes $2x^2-30x+63=0\implies x=\dfrac{15-3\sqrt{11}}{2}$

 Our answer is $8-(15-3\sqrt{11})=3\sqrt{11}-7\implies \boxed{\text{(B)}~7}$. (Mudkipswims42)

 



\textbf{Solution 2}

Denote the length of the equilateral octagon as $x$. The length of $\overline{BQ}$ can be expressed as $\frac{8-x}{2}$. By the Pythagorean Theorem, we find that:
 \[\left(\frac{8-x}{2}\right)^2+\overline{CQ}^2=x^2\implies \overline{CQ}=\sqrt{x^2-\left(\frac{8-x}{2}\right)^2}\] Since $\overline{CQ}=\overline{DR}$, we can say that $x+2\sqrt{x^2-\left(\frac{8-x}{2}\right)^2}=6\implies x=-7\pm3\sqrt{11}$. We can discard the negative solution, so $k+m+n=-7+3+11=\boxed{\textbf{(B) }7}$ ~ blitzkrieg21

 



\textbf{Solution 3}

Let the octagon's side length be $x$. Then $PH = \frac{6 - x}{2}$ and $PA = \frac{8 - x}{2}$. By the Pythagorean theorem, $PH^2 + PA^2 = HA^2$, so $\left(\frac{6 - x}{2} \right)^2 + \left(\frac{8 - x}{2} \right)^2 = x^2$. By expanding the left side and  combining the like terms, we get $\frac{x^2}{2} - 7x + 25 = x^2 \implies -\frac{x^2}{2} - 7x + 25 = 0$. Solving this using the quadratic formula, $\frac{-b \pm \sqrt{b^2-4ac}}{2a}$, we use $a = -\frac{1}{2}$, $b = -7$, and $c = 25$, to get one positive solution, $x=-7+3\sqrt{11}$, so $k+m+n=-7+3+11=\boxed{\textbf{(B) }7}$




\textbf{Solution 4}

Let $AB$, or the side of the octagon, be $x$. Then, $BQ = \left(\frac{8-x}{2}\right)$ and $CQ = \left(\frac{6-x}{2}\right)$. By the \href{https://artofproblemsolving.com/wiki/index.php?title=Pythagorean\_Theorem}{Pythagorean Theorem}, $BQ^2+CQ^2=x^2$, or $\left(\frac{8-x}{2}\right)^2+\left(\frac{6-x}{2}\right)^2 = x^2$. Multiplying this out, we have $x^2 = \frac{64-16x+x^2+36-12x+x^2}{4}$. Simplifying, $-2x^2-28x+100=0$. Dividing both sides by $-2$ gives $x^2+14x-50=0$. Therefore, using the \href{https://artofproblemsolving.com/wiki/index.php?title=Quadratic\_formula}{quadratic formula}, we have $x=-7 \pm 3\sqrt{11}$. Since lengths are always positive, then $x=-7+3\sqrt{11} \Rightarrow k+m+n=-7+3+11=\boxed{\textbf{(B)}\ 7}$



	\item (\textit{AMC 10 10B 2018}) Three young brother-sister pairs from different families need to take a trip in a van. These six children will occupy the second and third rows in the van, each of which has three seats. To avoid disruptions, siblings may not sit right next to each other in the same row, and no child may sit directly in front of his or her sibling. How many seating arrangements are possible for this trip?


 $\textbf{(A) }60 \qquad \textbf{(B) }72 \qquad \textbf{(C) }92 \qquad \textbf{(D) }96 \qquad \textbf{(E) }120 \qquad$

 \textbf{Answer: }D



\textbf{Solution 1 (Casework)}

We can begin to put this into cases. Let's call the pairs $a$, $b$ and $c$, and assume that a member of pair $a$ is sitting in the leftmost seat of the second row. We can have the following cases then.

 Case $1$: Second Row: a b cThird Row: b c a

 Case $2$:Second Row: a c bThird Row: c b a

 Case $3$: Second Row: a b cThird Row: c a b

 Case $4$: Second Row: a c bThird Row: b a c

 For each of the four cases, we can flip the siblings, as they are distinct. So, each of the cases has $2 \cdot 2 \cdot 2 = 8$ possibilities. Since there are four cases, when pair $a$ has someone in the leftmost seat of the second row, there are $32$ ways to rearrange it. However, someone from either pair $a$, $b$, or $c$ could be sitting in the leftmost seat of the second row. So, we have to multiply it by $3$ to get our answer of $32 \cdot 3 = 96$. So, the correct answer is 

 


	\item (\textit{AMC 10 10B 2018}) Joey and Chloe and their daughter Zoe all have the same birthday. Joey is $1$ year older than Chloe, and Zoe is exactly $1$ year old today. Today is the first of the $9$ birthdays on which Chloe's age will be an integral multiple of Zoe's age. What will be the sum of the two digits of Joey's age the next time his age is a multiple of Zoe's age?


 $\textbf{(A) }7 \qquad \textbf{(B) }8 \qquad \textbf{(C) }9 \qquad \textbf{(D) }10 \qquad \textbf{(E) }11 \qquad$

 \textbf{Answer: }E



\textbf{Solution 1}

Suppose that Chloe is $c$ years old today, so Joey is $c+1$ years old today. After $n$ years, Chloe and Zoe will be $n+c$ and $n+1$ years old, respectively. We are given that 
 \[\frac{n+c}{n+1}=1+\frac{c-1}{n+1}\] is an integer for $9$ nonnegative integers $n.$ It follows that $c-1$ has $9$ positive divisors. The prime factorization of $c-1$ is either $p^8$ or $p^2q^2.$ Since $c-1<100,$ the only possibility is $c-1=2^2\cdot3^2=36,$ from which $c=37.$ We conclude that Joey is $c+1=38$ years old today.

 Suppose that Joey's age is a multiple of Zoe's age after $k$ years, in which Joey and Zoe will be $k+38$ and $k+1$ years old, respectively. We are given that 
 \[\frac{k+38}{k+1}=1+\frac{37}{k+1}\] is an integer for some positive integer $k.$ It follows that $37$ is divisible by $k+1,$ so the only possibility is $k=36.$ We conclude that Joey will be $k+38=74$ years old then, from which the answer is $7+4=\boxed{\textbf{(E) }11}.$




\textbf{Solution 2}

Let Joey's age be $j$, Chloe's age be $c$, and we know that Zoe's age is $1$. 

 We know that there must be $9$ values $k\in\mathbb{Z}$ such that $c+k=a(1+k)$ where $a$ is an integer.

 Therefore, $c-1+(1+k)=a(1+k)$ and $c-1=(1+k)(a-1)$. Therefore, we know that, as there are $9$ solutions for $k$, there must be $9$ solutions for $c-1$. We know that this must be a perfect square. Testing perfect squares, we see that $c-1=36$, so $c=37$. Therefore, $j=38$. Now, since $j-1=37$, by similar logic, $37=(1+k)(a-1)$, so $k=36$ and Joey will be $38+36=74$ and the sum of the digits is $\boxed{\textbf{(E) }11}$.

 



\textbf{Solution 3}

Let Chloe's age be $c$ so that Zoey is 1. Then, if $n$ years pass, we have that again $c+n \equiv 0 \pmod{n+1}$ from which $c \equiv 1 \pmod{n+1}$ thus $c - 1 \equiv 0 \pmod{n+1}$. Hence, $c-1$ has 9 divisors. Assuming no human lives past a century, it is given that only numbers of the form $p^2q^2$ work. Taking minimality, $p = 2, q = 3$ from which $c-1 = 36 \implies c = 37$. Hence, Joey is 38. Now, if after some $k$ years Joey is a multiple of Zoe, we can just apply the same process as before. In other words, $38+k \equiv 0 \pmod{1+k}$ hence $38 \equiv 1 \pmod{1+k}$, and thus $37 \equiv 0 \pmod{1+k}$. The only integer multiples of 37 is 1 and 37, but they wish for the next integral multiple, which is when Zoe is 38 years and Joey is 74, whose sum is $\boxed{11}$.

 



\textbf{Solution 4}

Here's a different way of stating Solution 2:

 If a number is a multiple of both Chloe's age and Zoe's age, then it is a multiple of their difference. Since the difference between their ages does not change, then that means the difference between their ages has $9$ factors. Therefore, the difference between Chloe and Zoe's age is $36$, so Chloe is $37$, and Joey is $38$. The common factor that will divide both of their ages is $37$, so Joey will be $74$. The answer is $7 + 4 =  \boxed{\textbf{(E) }11}$.

 



\textbf{Solution 5}

Similar approach to above, just explained less concisely and more in terms of the problem (less algebraic).

 Let $C+n$ denote Chloe's age, $J+n$ denote Joey's age, and $Z+n$ denote Zoe's age, where $n$ is the number of years from now. We are told that $C+n$ is a multiple of $Z+n$ exactly nine times. Because $Z+n$ is $1$ at $n=0$ and will increase until greater than $C-Z$, it will hit every natural number less than $C-Z$, including every factor of $C-Z$. For $C+n$ to be an integral multiple of $Z+n$, the difference $C-Z$ must also be a multiple of $Z$, which happens if $Z$ is a factor of $C-Z$. Therefore, $C-Z$ has nine factors. The smallest number that has nine positive factors is $2^23^2=36$ . (We want it to be small so that Joey will not have reached three digits of age before his age is a multiple of Zoe's.) We also know $Z=1$ and $J=C+1$. Thus,
 \begin{align*} C-Z&=36, \\ J-Z&=37. \end{align*}
By our above logic, the next time $J-Z$ is a multiple of $Z+n$ will occur when $Z+n$ is a factor of $J-Z$. Because $37$ is prime, the next time this happens is at $Z+n=37$, when $J+n=74$. The answer is $7+4=\boxed{\textbf{(E) }11}$.

 



\textbf{Solution 6 (Number Theory+Logic)}

Let Joeys current age be $J$, Chloe's current age be $C$, Zoe's current age be $Z$, and $C-Z=C-1=K$. This gives $C=1+K$ and $J=2+K$. After $t$ years, Zoe will be $Z+t=t+1$ years old and Chloe will be $C+t$ years old. Chloe's age is an integral multiple of Zoe's age when $Z+t|C+t$. By the Divisibility of Differences rule (which is a unique case of the linear combination property of divisibility), $t+1|C+t-(t+1)$->$t+1|C-1$->$t+1=K$.

 $K$ must have exactly $9$ positive integer factors, since there are precisely $9$ values of $t$ for which Zoe's age is a factor of Chloe's age. The smallest number with this property is $2^2 \times 3^2=36$. We assume this is the value of $K$, as the next smallest value of $K$ (which is $K=100$) yields the unlikely case that Chloe is currently $101$ years old. 

 Since $K=36$, we have $C=1+36=37$ and $J=2+36=38$. Joey's age is a multiple of Zoe's age when $Z+t|J+t$. Plugging in the values of $Z$ and $J$, we get $t+1|t+38$. By the Divisibility of Differences rule, $t+1|t+38-(t+1)=37$. The smallest value of $t$ is $36$, so Joey will be $38+36=74$ years old. The sum of the digits is $7+4=\boxed{\textbf{(E) } 11}$.

 


	\item (\textit{AMC 10 10B 2018}) A function $f$ is defined recursively by $f(1)=f(2)=1$ and 
 \[f(n)=f(n-1)-f(n-2)+n\]for all integers $n \geq 3$. What is $f(2018)$?


 $\textbf{(A) } 2016 \qquad \textbf{(B) } 2017 \qquad \textbf{(C) } 2018 \qquad \textbf{(D) } 2019 \qquad \textbf{(E) } 2020$

 \textbf{Answer: }B



\textbf{Solution 1 (Algebra)}

For all integers $n \geq 7,$ note that
 \begin{align*} f(n)&=f(n-1)-f(n-2)+n \\ &=[f(n-2)-f(n-3)+n-1]-f(n-2)+n \\ &=-f(n-3)+2n-1 \\ &=-[f(n-4)-f(n-5)+n-3]+2n-1 \\ &=-f(n-4)+f(n-5)+n+2 \\ &=-[f(n-5)-f(n-6)+n-4]+f(n-5)+n+2 \\ &=f(n-6)+6. \end{align*}
It follows that
 \begin{align*} f(2018)&=f(2012)+6 \\ &=f(2006)+12 \\ &=f(2000)+18 \\ & \ \vdots \\ &=f(2018-6\cdot336)+6\cdot336 \\ &=f(2)+2016 \\ &=\boxed{\textbf{(B) } 2017}. \end{align*}
~MRENTHUSIASM

 



\textbf{Solution 2 (Algebra)}

For all integers $n\geq3,$ we rearrange the given equation: 
 \[f(n)-f(n-1)+f(n-2)=n. \hspace{28.25mm}(1)\]For all integers $n\geq4,$ it follows that 
 \[f(n-1)-f(n-2)+f(n-3)=n-1. \hspace{15mm}(2)\]For all integers $n\geq4,$ we add $(1)$ and $(2):$ 
 \[f(n)+f(n-3)=2n-1. \hspace{38.625mm}(3)\]For all integers $n\geq7,$ it follows that 
 \[f(n-3)+f(n-6)=2n-7. \hspace{32mm}(4)\]For all integers $n\geq7,$ we subtract $(4)$ from $(3):$ 
 \[f(n)-f(n-6)=6. \hspace{47.5mm}(5)\]From $(5),$ we have the following system of $\frac{2018-8}{6}+1=336$ equations:
 \begin{align*} f(2018)-f(2012)&=6, \\ f(2012)-f(2006)&=6, \\ f(2006)-f(2000)&=6, \\ & \ \vdots \\ f(8)-f(2)&=6. \end{align*}
We add these equations up to get 
 \[f(2018)-f(2)=6\cdot336=2016,\] from which $f(2018)=f(2)+2016=\boxed{\textbf{(B) } 2017}.$




\textbf{Solution 3 (Finite Differences)}

Preamble: In this solution, we define the sequence $A$ to satisfy $a_n = f(n),$ where $a_n$ represents the $n$th term of the sequence $A.$ This solution will show a few different perspectives. Even though it may not be as quick as some of the solutions above, I feel like it is an interesting concept, and may be more motivated. 

 To begin, we consider the sequence $B$ formed when we take the difference of consecutive terms between $A.$ Define $b_n = a_{n+1} - a_n.$ Notice that for $n \ge 4,$ we have 

 $\begin{cases}\begin{aligned} a_{n+1} &= a_{n} - a_{n-1} + (n+1) \\ a_{n} &= a_{n-1} - a_{n-2} + n \end{aligned}.\end{cases}$

  Notice that subtracting the second equation from the first, we see that $b_{n} = b_{n-1} - b_{n-2} + 1.$ 

 If you didn’t notice that $B$ repeated directly in the solution above, you could also, possibly more naturally, take the finite differences of the sequence $b_n$ so that you could define $c_n = b_{n+1} - b_n.$ Using a similar method as above through reindexing and then subtracting, you could find that $c_n = c_{n-1} - c_{n-2}.$ The sum of any six consecutive terms of a sequence which satisfies such a recursion is $0,$ in which you have that $b_{n} = b_{n+6}.$ In the case in which finite differences didn’t reduce to such a special recursion, you could still find the first few terms of $C$ to see if there are any patterns, now that you have a much simpler sequence. Doing so in this case, it can also be seen by seeing that the sequence $C$ looks like 
 \[\underbrace{2, 1, -1, -2, -1, 1,}_{\text{cycle period}} 2, 1, -1, -2, -1, 1, \ldots\] in which the same result follows. 

 Using the fact that $B$ repeats every six terms, this motivates us to look at the sequence $B$ more carefully. Doing so, we see that $B$ looks like 
 \[\underbrace{2, 3, 2, 0, -1, 0,}_{\text{cycle period}} 2, 3, 2, 0, -1, 0, \ldots\] (If you tried pattern finding on sequence $B$ directly, you could also arrive at this result, although I figured defining a second sequence based on finite differences was more motivated.)

 Now, there are two ways to finish. 

 Finish Method \#1: Notice that any six consecutive terms of $B$ sum to $6,$ after which we see that $a_n = a_{n-6} + 6.$ Therefore, $a_{2018} = a_{2012} + 6 = \cdots = a_{2} + 2016 = \boxed{\textbf{(B) } 2017}.$

 Finish Method \#2: Notice that $a_{2018} = a_{2017} - a_{2016} + 2018 = B_{2016} + 2018 = -1 + 2018 = \boxed{\textbf{(B) } 2017}.$




\textbf{Solution 4 (Pattern)}

Start out by listing some terms of the sequence. 
 \begin{align*} f(1)&=1 \\ f(2)&=1 \\ f(3)&=3 \\ f(4)&=6 \\ f(5)&=8 \\ f(6)&=8 \\  f(7)&=7 \\ f(8)&=7 \\ f(9)&=9 \\ f(10)&=12 \\ f(11)&=14 \\ f(12)&=14 \\ f(13)&=13 \\ f(14)&=13 \\ f(15)&=15 \\ & \ \vdots \end{align*}
Notice that $f(n)=n$ whenever $n$ is an odd multiple of $3$, and the pattern of numbers that follow will always be $+3$, $+2$, $+0$, $-1$, $+0$, $+2$.The largest odd multiple of $3$ smaller than $2018$ is $2013$, so we have
 \begin{align*} f(2013)&=2013 \\ f(2014)&=2016 \\ f(2015)&=2018 \\ f(2016)&=2018 \\ f(2017)&=2017 \\ f(2018)&=\boxed{\textbf{(B) } 2017}. \end{align*}





\textbf{Solution 5 (Pattern)}

Writing out the first few values, we get
 \[1,1,3,6,8,8,7,7,9,12,14,14,13,13,15,18,20,20,19,19,\ldots.\] We see that every number $x$ where $x \equiv 1\pmod 6$ has $f(x)=x,f(x+1)=f(x)=x,$ and $f(x-1)=f(x-2)=x+1.$ The greatest number that's $1\pmod{6}$ and less than $2018$ is $2017,$ so we have $f(2017)=f(2018)=\boxed{\textbf{(B) } 2017}.$




\textbf{Solution 6 (Complex Numbers)}


 \begin{align*} f(n)&=f(n-1)-f(n-2)+n \\ f(n-1)&=f(n-2)-f(n-3)+n-1 \end{align*}
Subtracting those two and rearranging gives 
 \begin{align*} f(n)-2f(n-1)+2f(n-2)-f(n-3)&=1 \\ f(n-1)-2f(n-2)+2f(n-3)-f(n-4)&=1 \end{align*}
Subtracting those two gives $f(n)-3f(n-1)+4f(n-2)-3f(n-3)+f(n-4)=0.$

 The characteristic polynomial is $x^4-3x^3+4x^2-3x+1=0.$
$x=1$ is a root, so using synthetic division results in $(x-1)(x^3-2x^2+2x-1)=0.$
$x=1$ is a root, so using synthetic division results in $(x-1)^2(x^2-x+1)=0.$
$x^2-x+1=0$ has roots $x=\frac{1}{2}\pm\frac{i\sqrt{3}}{2}.$

 And 
 \[f(n)=(An+D)\cdot1^n+B\cdot\left(\frac{1}{2}-\frac{i\sqrt{3}}{2}\right)^n+C\cdot\left(\frac{1}{2}+\frac{i\sqrt{3}}{2}\right)^n.\]Plugging in $n=1$, $n=2$, $n=3$, and $n=4$ results in a system of $4$ linear equations$\newline$Solving them gives $A=1, \ B=\frac{1}{2}-\frac{i\sqrt{3}}{2}, \ C=\frac{1}{2}+\frac{i\sqrt{3}}{2}, \ D=1.$ Note that you can guess $A=1$ by answer choices.

 So plugging in $n=2018$ results in 
 \begin{align*} 2018+1+\left(\frac{1}{2}-\frac{i\sqrt{3}}{2}\right)^{2019}+\left(\frac{1}{2}+\frac{i\sqrt{3}}{2}\right)^{2019}&=2019+(\cos(-60^{\circ})+i\sin(-60^{\circ}))^{2019})+(\cos(60^{\circ})+i\sin(60^{\circ}))^{2019}) \\ &=2019+(\cos(-60^{\circ}\cdot2019)+i\sin(-60^{\circ}\cdot2019))+(\cos(60^{\circ}\cdot2019)+i\sin(60^{\circ}\cdot2019)) \\ &=\boxed{\textbf{(B) } 2017}. \end{align*}
~ryanbear

 



\textbf{Solution 7}

We utilize patterns to solve this equation:
 \begin{align*} f(3)&=3, \\ f(4)&=6, \\ f(5)&=8, \\ f(6)&=8, \\ f(7)&=7, \\ f(8)&=7. \end{align*}
We realize that the pattern repeats itself. For every six terms, there will be four terms that we repeat, and two terms that we don't repeat. We will exclude the first two for now, because they don't follow this pattern. 

 First, we need to know whether or not $2016$ is part of the skip or repeat. We notice that $f(6),f(12), \ldots,f(6n)$ all satisfy $6+6(n-1)=n,$ and we know that $2016$ satisfies this, leaving $n=336.$ Therefore, we know that $2016$ is part of the repeat section. But what number does it repeat? 

 We know that the repeat period is $2,$ and it follows that pattern of $1,1,8,8,7,7.$ Again, since $f(6) = f(5)$ and so on  for the repeat section, $f(2016)=f(2015),$ so we don't need to worry about which one, since it repeats with period $2.$ We see that the repeat pattern of $f(6),f(12),\ldots,f(6n)$ follows $8,14,20,$ it is an arithmetic sequence with common difference $6.$ Therefore, $2016$ is the $335$th term of this, but including $1,$ it is $336\cdot6+1=\boxed{\textbf{(B) } 2017}.$

 
Edited by TheBinaryBlackHole

 



\textbf{Solution 8}

We list out the first $13$ terms and their differences: 

 
                    \begin{figure}[H]
                        \centering
                        \includegraphics[width=0.5\textwidth]{Solution/AMC_10_10B_2018_20_Solution_Figure_1.png}
                    \end{figure}
                     We see the differences are clearly repeating every $6$ terms here; we see that $2018\equiv2\pmod6$. We see that the values for $2$ and $8$ are both one less, so we conclude our answer is $2018-1=\boxed{\textbf{(B)}~2017}.$



	\item (\textit{AMC 10 10B 2018}) Mary chose an even $4$-digit number $n$. She wrote down all the divisors of $n$ in increasing order from left to right: $1,2,\ldots,\dfrac{n}{2},n$. At some moment Mary wrote $323$ as a divisor of $n$. What is the smallest possible value of the next divisor written to the right of $323$?


 $\textbf{(A) }324 \qquad \textbf{(B) }330 \qquad \textbf{(C) }340 \qquad \textbf{(D) }361 \qquad \textbf{(E) }646 \qquad$

 \textbf{Answer: }C



\textbf{Solution 1 (Inequalities)}

Let $d$ be the next divisor written to the right of $323.$

 If $\gcd(323,d)=1,$ then 
 \[n\geq323d>323^2>100^2=10000,\] which contradicts the precondition that $n$ is a $4$-digit number.

 It follows that $\gcd(323,d)>1.$ Since $323=17\cdot19,$ the smallest possible value of $d$ is $17\cdot20=\boxed{\textbf{(C) } 340},$ from which 
 \[n=\operatorname{lcm}(323,d)=17\cdot19\cdot20=6460.\]~MRENTHUSIASM ~tdeng

 



\textbf{Solution 2 (Inequalities)}

Let $d$ be the next divisor written to the right of $323.$

 Since $n$ is even and $323=17\cdot19,$ we have $n=2\cdot17\cdot19\cdot k=646k$ for some positive integer $k.$ Moreover, since $1000\leq n\leq9998,$ we get $2\leq k\leq15.$ As $d>323,$ it is clear that $d$ must be divisible by $17$ or $19$ or both.

 Therefore, the smallest possible value of $d$ is $17\cdot20=\boxed{\textbf{(C) } 340},$ from which 
 \[n=\operatorname{lcm}(323,d)=17\cdot19\cdot20=6460.\]~MRENTHUSIASM ~bjhhar

 



\textbf{Solution 3 (Quick)}

The prime factorization of $323$ is $17 \cdot 19$. Our answer must be a multiple of either $17$ or $19$ or both. Since $17 < 19$, the next smallest divisor that is divisble by $17$ would be $323 + 17 = \boxed{\textbf{(C) } 340}$




\textbf{Solution 4 (Answer Choices)}

Since prime factorizing $323$ gives you $17 \cdot 19$, the desired answer needs to be a multiple of $17$ or $19$, this is because if it is not a multiple of $17$ or $19$, $n$ will be more than a $4$ digit number. For example, if the answer were to instead be $324$, $n$ would have to be a multiple of $2^2\cdot3^4\cdot17\cdot19$ for both $323$ and $324$ to be a valid factor, meaning $n$ would have to be at least $104652$, which is too big. Looking at the answer choices, $\textbf{(A)}$ and $\textbf{(B)}$ are both not a multiple of neither $17$ nor $19$, $\textbf{(C)}$ is divisible by $17$. $\textbf{(D)}$ is divisible by $19$, and $\textbf{(E)}$ is divisible by both $17$ and $19$. Since $\boxed{\textbf{(C) } 340}$ is the smallest number divisible by either $17$ or $19$ it is the answer. Checking, we can see that $n$ would be $6460$, a $4$-digit number. Note that $n$ is also divisible by $2$, one of the listed divisors of $n$. (If $n$ was not divisible by $2$, we would need to look for a different divisor.)

 -Edited by Mathandski

 



\textbf{Solution 5 (Answer Choices)}

Note that $323$ multiplied by any of the answer choices results in a $5$ or $6$-digit $n$. So, we need a choice that shares a factor(s) with $323$, such that the factors we'll need to add to the prime factorization of $n$ (in result to adding the chosen divisor) won't cause our number to multiply to more than $4$ digits. The prime factorization of $323$ is $17\cdot19$, and since we know $n$ is even, our answer needs to be 

 
            \begin{itemize}
        \item even
            \end{itemize}
            
            \begin{itemize}
        \item has a factor of $17$ or $19$
            \end{itemize}
            We see $340$ achieves this and is the smallest to do so ($646$ being the other). So, we get $\boxed{\textbf{(C) } 340}$.

 



\textbf{Solution 6 (Olympiad Proof )}

This is not the fastest solution, but if i saw this question on an Olympiad/AIME, where there are no answer choices, and my work counted, this is what i would do (for the purpose of this question, a=323, b=answer, c=the four digit number);

 $\newline$Note that for any $a,b,c \in \mathbb{Z}^+$, if $a|c$ and $b|c$ then lcm$(a,b)|c$ (this is because if something divides a or b, it must divide c, and thus the max of all of the prime factors, ie: lcm, divides c) since $a=323,b > a$, if $gcd(a,b)=1$ then $lcm(a,b)=ab$ and thus $ab|c \implies c>a^2 \implies c>100^2 \implies$ c is not a four digit number.$\newline$ thus, $gcd(a,b)\neq1$. This implies that either $17|b$, or $19|b$, or both.

 $\newline$ Case 1: $17|b$, $19\not|b$. We let $b=17b'$, and by Euclid's Lemma, $19\not|b'$. Then, $lcm(323,b)|c \implies 17(lcm(19,b'))|c$. Since we already established that, $19\not|b$ (and since 19 is prime, if it does not divide a number it is coprime to that number), $17*19*b'|c \implies 323b'|c$. Since $b=17b'>19*17$, $b' \geq 20$. A quick check shows $b'=20, b=340$ suffices.$\newline$ Now, let us show that there are no such numbers less than 340.$\newline$ Presume there exists such a number, $n \in \mathbb{Z}^+$ is in the range $(323,340)$. By hypothesis, there is a $d>1$ such that $d|323$, $d|n$. By properties of divisibility $d|n-323$. the maximum possible value of $n-323$ is $17^-$ (basically an arbitrary amount smaller than 17). But, since $d>1$ and $d|323, d \in \{17,19,323\}$. Of which, the minimum value is d=17. but, $17>17^-$ so there is no such d, and no such n.$\newline$ Thus, our answer is just $\boxed{\textbf{(C) } 340}$.$\newline$ ~Stereotypicalmathnerd

 



\textbf{Solution 7}

Let the value of the next divisor be $x$. Then it follows that $lcm(323,x)|n$. For the sake of this problem, suppose $lcm(323,x)=n$, since it seems unlikely that $lcm(323,x)\leq 4999$. From here, there are two ways to solve this problem (note that $323=17\times19)$. $\newline$
$\textbf{Ending 1.1 }$ 

 You can simply prime factorize each answer choice and find one that will produce a valid value for $n$. $\newline$
$\textbf{Ending 1.2 }$ 

 Alternatively, if you know the identity relating $gcd$ and $lcm$ ($gcd(a,b)\times lcm(a,b)=a\times b$), this problem can be rewritten as $gcd(323,x)\times lcm(323,x)=323\times x$. Since $lcm(323,x)=n$, this means $n=\frac{323\times x}{gcd(323,x)}$.

 Using either ending will produce the final answer of $\boxed{\textbf{(C) } 340}$ (This means $n=\frac{323\times 340}{gcd(323,340)} \to n=\frac{17^2 \times 19 \times 20}{17} \to n=6460$, which checks out).$\newline$




\textbf{Solution 8}

The key point is that $n$ is even. Also, $323n < 10000 \implies n \leq 30$. The key idea is that $323 = 17 \cdot 19$. If we were to divide by some $n < 17$, then dividing by 17 yields a smaller value, which is not allowed. Hence, consider $n = 18$ in which $19 \cdot 18 = 342$ is obtained. 

 Now we check the next value, when $n = 20$ so we can check 19. In that case, $17 \cdot 20 = 340$ is formed.

 We now prove no other values exist. Simply speaking, if $n > 20$, then consider $323 \cdot n$. If $n$ is a prime, dividing by $n$ gives 323, and the next number is anyways $17 \cdot n$. This however is minimized at $n = 20$, hence no cases arise. Similarly, if $n$ is composite with a prime factorization, one can still rewrite it of the form $323 \cdot n$. But all prime factors of $n$ are less than or equal to 15, hence the minimum is still at when we divide by 19, in which $17 \cdot n$. But this is still minimized at $n = 20$, hence no cases arise. $\blacksquare$

 The answer is $\boxed{340}$, and we are done. $\square$



	\item (\textit{AMC 10 10B 2018}) Real numbers $x$ and $y$ are chosen independently and uniformly at random from the interval $[0,1]$. Which of the following numbers is closest to the probability that $x,y,$ and $1$ are the side lengths of an obtuse triangle?


 $\textbf{(A) }0.21 \qquad \textbf{(B) }0.25 \qquad \textbf{(C) }0.29 \qquad \textbf{(D) }0.50 \qquad \textbf{(E) }0.79 \qquad$

 \textbf{Answer: }C



\textbf{Solution 1}

The Pythagorean Inequality tells us that in an obtuse triangle, $a^{2} + b^{2} < c^{2}$. The triangle inequality tells us that $a + b > c$. So, we have two inequalities:
 \[x^2 + y^2 < 1\]
 \[x + y > 1\]The first equation is $\frac14$ of a circle with radius $1$, and the second equation is a line from $(0, 1)$ to $(1, 0)$.So, the area is $\frac{\pi}{4} - \frac12$ which is approximately $\boxed{\textbf{(C)} ~0.29}$




\textbf{Solution 2 (Trig)}

Note that the obtuse angle in the triangle has to be opposite the side that is always length $1$. This is because the largest angle is always opposite the largest side, and if two sides of the triangle were $1$, the last side would have to be greater than $1$ to make an obtuse triangle. Using this observation, we can set up a law of cosines where the angle is opposite $1$:

 
 \[1^2=x^2+y^2-2xy\cos(\theta)\]

 where $x$ and $y$ are the sides that go from $[0,1]$ and $\theta$ is the angle opposite the side of length $1$.

 By isolating $\cos(\theta)$, we get:

 
 \[\frac{1-x^2-y^2}{-2xy} = \cos(\theta)\]

 For $\theta$ to be obtuse, $\cos(\theta)$ must be negative. Therefore, $\frac{1-x^2-y^2}{-2xy}$ is negative. Since $x$ and $y$ must be positive, $-2xy$ must be negative, so we must make $1-x^2-y^2$ positive. From here, we can set up the inequality
 \[x^2+y^2<1\]Additionally, to satisfy the definition of a triangle, we need:
 \[x+y>1\]The solution should be the overlap between the two equations in the first quadrant.

 By observing that $x^2+y^2<1$ is the equation for a circle, the amount that is in the first quadrant is $\frac{\pi}{4}$. The line can also be seen as a chord that goes from $(0, 1)$ to $(1, 0)$. By cutting off the triangle of area $\frac{1}{2}$ that is not part of the overlap, we get $\frac{\pi}{4} - \frac{1}{2} \approx \boxed{\textbf{(C)} ~0.29}$.

 ..why would you do this? for what purpose? its much more complicated and this is the AMC 10! -Orion 2010

 


	\item (\textit{AMC 10 10B 2018}) How many ordered pairs $(a, b)$ of positive integers satisfy the equation 

 \[a\cdot b + 63 = 20\cdot \text{lcm}(a, b) + 12\cdot\text{gcd}(a,b),\]
where $\text{gcd}(a,b)$ denotes the greatest common divisor of $a$ and $b$, and $\text{lcm}(a,b)$ denotes their least common multiple?


 $\textbf{(A)} \text{ 0} \qquad \textbf{(B)} \text{ 2} \qquad \textbf{(C)} \text{ 4} \qquad \textbf{(D)} \text{ 6} \qquad \textbf{(E)} \text{ 8}$

 \textbf{Answer: }B



\textbf{Solution}

Let $x = \text{lcm}(a, b)$, and $y = \text{gcd}(a, b)$. Therefore, $a\cdot b = \text{lcm}(a, b)\cdot \text{gcd}(a, b) = x\cdot y$. Thus, the equation becomes

 
 \[x\cdot y + 63 = 20x + 12y\]
 \[x\cdot y - 20x - 12y + 63 = 0\]

 Using \href{https://artofproblemsolving.com/wiki/index.php?title=Simon\%27s\_Favorite\_Factoring\_Trick}{Simon's Favorite Factoring Trick}, we rewrite this equation as

 
 \[(x - 12)(y - 20) - 240 + 63 = 0\]
 \[(x - 12)(y - 20) = 177\]

 Since $177 = 3\cdot 59$ and $x >= y$, we have $x  - 12 = 59$ and $y - 20 = 3$, or $x - 12 = 177$ and $y - 20 = 1$. This gives us the solutions $(71, 23)$ and $(189, 21)$. Since the $\text{GCD}$ must be a divisor of the $\text{LCM}$, the first pair does not work. Assume $a>b$. We must have $a = 21 \cdot 9$ and $b = 21$, and we could then have $a<b$, so there are $\boxed{\textbf{(B)} ~2}$ solutions.(awesomeag)

 Edited by IronicNinja, Firebolt360, and mprincess0229~

 


	\item (\textit{AMC 10 10B 2018}) Let $ABCDEF$ be a regular hexagon with side length $1$. Denote by $X$, $Y$, and $Z$ the midpoints of sides $\overline {AB}$, $\overline{CD}$, and $\overline{EF}$, respectively. What is the area of the convex hexagon whose interior is the intersection of the interiors of $\triangle ACE$ and $\triangle XYZ$?


 $\textbf{(A)}\ \frac {3}{8}\sqrt{3} \qquad \textbf{(B)}\ \frac {7}{16}\sqrt{3} \qquad \textbf{(C)}\ \frac {15}{32}\sqrt{3} \qquad  \textbf{(D)}\ \frac {1}{2}\sqrt{3} \qquad \textbf{(E)}\ \frac {9}{16}\sqrt{3}$

 \textbf{Answer: }C



\textbf{Solution 1 (Area Addition)}


                    \begin{figure}[H]
                        \centering
                        \includegraphics[width=0.5\textwidth]{Solution/AMC_10_10B_2018_24_Solution_Figure_1.png}
                    \end{figure}
                    

 The desired area (hexagon $MPNQOR$) consists of an equilateral triangle ($\triangle MNO$) and three right triangles ($\triangle MPN,\triangle NQO,$ and $\triangle ORM$).

 Notice that $\overline {AD}$ (not shown) and $\overline {BC}$ are parallel. $\overline {XY}$ divides transversals $\overline {AB}$ and $\overline {CD}$ into a $1:1$ ratio (This can be shown by similar triangles.). Thus, it must also divide transversal $\overline {AC}$ and transversal $\overline {CO}$ into a $1:1$ ratio. By symmetry, the same applies for $\overline {CE}$ and $\overline {EA}$ as well as $\overline {EM}$ and $\overline {AN}.$

 In $\triangle ACE,$ we see that $\frac{[MNO]}{[ACE]} = \frac{1}{4}$ and $\frac{[MPN]}{[ACE]} = \frac{1}{8}.$ Our desired area becomes 
 \[\left(\frac{1}{4}+3 \cdot \frac{1}{8}\right) \cdot \frac{(\sqrt{3})^2 \cdot \sqrt{3}}{4} = \boxed{\textbf{(C)}\ \frac {15}{32}\sqrt{3}}.\]




\textbf{Solution 2 (Area Subtraction)}

Instead of directly finding the desired hexagonal area, $\triangle XYZ$ can be found. It consists of three triangles and the desired hexagon. Given triangle rotational symmetry, the three triangles are congruent. See that $\triangle XYZ$ and $\triangle ACE$ are equilateral, so $m\angle PXN=60,$ so $m\angle AXP = \frac{180-60}{2}=60.$ As $\overline {AC}$ is a transversal running through $\overline {FC}$ (use your imagination) and $\overline {AB},$ we have $m\angle BAC=m\angle FCA = \frac{m\angle ACE}{2}=30.$ 

 Then, $\triangle APX$ is a $30$-$60$-$90$ triangle. By HL congruence, $\triangle APX \cong \triangle NPX.$ Note that $AX=\frac{1}{2}.$ Then, the area of $\triangle PXN$ is $\frac{\sqrt{3}}{32}.$ There are three such triangles for a total area of $\triangle XYZ$ is $\frac{3\sqrt{3}}{32}.$ Find the side of $\triangle XYZ$ to be $\frac{3}{2},$ so the area is $\frac{9\sqrt{3}}{16}.$ 

 The answer is 
 \[\frac{9\sqrt{3}}{16}-\frac{3\sqrt{3}}{32}=\boxed{\textbf{(C)}\ \frac {15}{32}\sqrt{3}}.\]~BJHHar

 



\textbf{Solution 3 (Area Subtraction)}

We will find its area by subtracting the three triangles congruent to triangle APM from $\triangle ACE$. By the rotational symmetry of the regular hexagon, the three corner triangles outside the intersection ($\triangle APM$, $\triangle CQN$, and $\triangle ERO$) are all congruent. 

 We first find the area of the large equilateral triangle $\triangle ACE$. Since the regular hexagon has a side length of $1$, the side length of $\triangle ACE$ is $\sqrt{3}.$ Using the area formula for an equilateral triangle, we have:
 \[[\text{ACE}] = \frac{(\sqrt{3})^2\sqrt{3}}{4} = \frac{3\sqrt{3}}{4}.\]

 To find the area of the corner triangle $\triangle APM$, we can compare it to the smaller adjacent triangle $\triangle ZMR.$ Since $\triangle APM$ and $\triangle ZMR$ share the same vertical angles at vertex $M$ and are bounded by the grid lines of the hexagon, analyzing their bases and heights reveals that the area of $\triangle APM$ is exactly $3$ times the area of $\triangle ZMR.$ Thus, the area ratio between $\triangle APM$ and $\triangle ZMR$ is $3:1.$ 

 The area of the smaller triangle $\triangle ZMR$ can be calculated from the grid as:
 \[[\text{ZMR}] = \frac{1}{4} \cdot \frac{\sqrt{3}}{4} \cdot \frac{1}{2} = \frac{\sqrt{3}}{32}.\]

 Multiplying this by $3$ gives the true area of one corner triangle $\triangle APM$:
 \[[\text{APM}] = 3 \cdot \frac{\sqrt{3}}{32} = \frac{3\sqrt{3}}{32}.\]

 Since there are three such identical corner triangles surrounding the inner hexagon, we subtract their total area from the area of $\triangle ACE$:

 The answer is 
 \[\frac{3\sqrt{3}}{4} - 3\left(\frac{3\sqrt{3}}{32}\right) = \frac{24\sqrt{3}}{32} - \frac{9\sqrt{3}}{32} = \boxed{\textbf{(C)}\ \frac {15}{32}\sqrt{3}}.\]




\textbf{Solution 4 (Area Subtraction)}


                    \begin{figure}[H]
                        \centering
                        \includegraphics[width=0.5\textwidth]{Solution/AMC_10_10B_2018_24_Solution_Figure_1.png}
                    \end{figure}
                    

 Now, if we look at the figure, we can see that the complement of the hexagon we are trying to find is composed of $3$ isosceles trapezoids (namely $AXZF,CYXB,$ and $EZYD$) and $3$ right triangles (namely $\triangle XPN,\triangle YQO,$ and $\triangle ZRM$).

 Finding the trapezoid's area, we know that one base of each trapezoid is just the side length of the hexagon, which is $1,$ and the other base is $\frac{3}{2}$ (it is halfway in between the side and the longest diagonal, which has length $2$) with a height of $\frac{\sqrt{3}}{4}$ (by using the Pythagorean Theorem and the fact that it is an isosceles trapezoid) to give each trapezoid having an area of $\frac{5\sqrt{3}}{16}$ for a total area of $\frac{15\sqrt{3}}{16}.$ (Alternatively, we could have calculated the area of hexagon $ABCDEF$ and subtracted the area of $\triangle XYZ,$ which, as we showed before, had a side length of $\frac{3}{2}$).

 Now, we need to find the area of each of the small triangles, which, if we look at the triangle that has a vertex on $X,$ is similar to the triangle with a base of $YC = \frac12.$ Using similar triangles, we calculate the base to be $\frac{1}{4}$ and the height to be $\frac{\sqrt{3}}{4}$ giving us an area of  $\frac{\sqrt{3}}{32}$ per triangle, and a total area of $\frac{3\sqrt{3}}{32}.$ Adding the two areas together, we get $\frac{15\sqrt{3}}{16} + \frac{3\sqrt{3}}{32} = \frac{33\sqrt{3}}{32}.$ Finding the total area, we get $6 \cdot 1^2 \cdot \frac{\sqrt{3}}{4}=\frac{3\sqrt{3}}{2}.$ Taking the complement, we get $\frac{3\sqrt{3}}{2} - \frac{33\sqrt{3}}{32} = \boxed{\textbf{(C)}\ \frac {15}{32}\sqrt{3}}.$




\textbf{Solution 5 (Area Subtraction)}


                    \begin{figure}[H]
                        \centering
                        \includegraphics[width=0.5\textwidth]{Solution/AMC_10_10B_2018_24_Solution_Figure_3.png}
                    \end{figure}
                    

 We could also subtract $\triangle APM,\triangle CQN,$ and $\triangle ERO$ from $\triangle ACE.$

 Since $\angle BAF = 120^{\circ}$ and $\angle BAC = \angle FAE = 30^{\circ},$ we have $\angle CAE = \angle BAF-\angle BAC-\angle FAE=60^{\circ}.$

 Since $AX=BX$ and $FZ=EZ,$ we have $AF \parallel XZ,$ from which $\angle AMX= \angle FAM = 30^{\circ}.$

 We can show that $\triangle APM$ is $30$-$60$-$90$ using a similar method, $\triangle CQN$ and $\triangle ERO$ are also $30$-$60$-$90.$

 Since $AC=CE=AE=\sqrt{3},$ we have $[ACE]=AC^2 \cdot \frac{\sqrt{3}}{4}=3 \cdot \frac{\sqrt{3}}{4} = \frac{3 \sqrt{3}}{4}.$

 Since $AX= \frac{1}{2}$ and $AP = AX \cdot \frac{\sqrt{3}}{2}= \frac{1}{2} \cdot \frac{\sqrt{3}}{2} = \frac{\sqrt{3}}{4},$ we have $PM = AP \cdot \sqrt{3} = \frac{\sqrt{3}}{4} \cdot \sqrt{3} = \frac{3}{4}.$

 Note that 
 \[[APM]=[CQN]=[ERO]=\frac{1}{2} \cdot AP \cdot PM = \frac{1}{2} \cdot \frac{\sqrt{3}}{4} \cdot \frac{3}{4} = \frac{3 \sqrt{3}}{32}.\]Therefore, we get
 \[[PNQORM]=[ACE]-[APM]-[CQN]-[ERO]=\frac{3 \sqrt{3}}{4} - 3 \cdot \frac{3 \sqrt{3}}{32} = \frac{24 \sqrt{3}}{32} - \frac{9 \sqrt{3}}{32} = \boxed{\textbf{(C)}\ \frac {15}{32}\sqrt{3}}.\]




\textbf{Solution 6 (Partition the Hexagon)}

We partition hexagon $ABCDEF$ into $48$ congruent $30^\circ\text{-}60^\circ\text{-}90^\circ$ triangles, as shown below:
                    \begin{figure}[H]
                        \centering
                        \includegraphics[width=0.5\textwidth]{Solution/AMC_10_10B_2018_24_Solution_Figure_4.png}
                    \end{figure}
                    Let the brackets denote areas. Note that the desired region contains $15$ of the $48$ small triangles, so the answer is 
 \[\frac{15}{48}[ABCDEF]=\frac{15}{48}\cdot\frac{3\sqrt3}{2}=\boxed{\textbf{(C)}\ \frac {15}{32}\sqrt{3}}.\]~AlexLikeMath ~MRENTHUSIASM

 



\textbf{Solution 7 (Partition the Hexagon)}

Dividing $\triangle MNO$ into two right triangles congruent to $\triangle PMN,$ we see that $[MPNQOR]=\dfrac{5}{8}[ACE].$ Because $[ACE] = \dfrac{1}{2}[ABCDEF],$ we have $[MPNQOR]=\dfrac{5}{16}[ABCDEF].$ From here, you should be able to tell that the answer will have a factor of $5,$ and $\boxed{\textbf{(C)}\ \frac {15}{32}\sqrt{3}}$ is the only answer that has a factor of $5.$ However, if you want to actually calculate the area, you would calculate $[ABCDEF]$ to be $6 \cdot \dfrac{\sqrt{3}}{2 \cdot 2} = \dfrac{3\sqrt{3}}{2},$ so $[MPNQOR] = \dfrac{5}{16} \cdot \dfrac{3\sqrt{3}}{2} = \frac {15}{32}\sqrt{3}.$




\textbf{Solution 8 (Trigonometry)}

Notice, the area of the convex hexagon formed through the intersection of the $2$ triangles can be found by finding the area of the triangle formed by the midpoints of the sides and subtracting the smaller triangles that are formed by the region inside this triangle but outside the other triangle. First, let's find the area of the triangle formed by the midpoint of the sides. Notice, this is an equilateral triangle, thus all we need is to find the length of its side. To do this, we look at the isosceles trapezoid outside this triangle but inside the outer hexagon. Since the interior angle of a regular hexagon is $120^\circ$ and the trapezoid is isosceles, we know that the angle opposite is $60^\circ,$ and thus the side length of this triangle is $1+2\left(\frac{1}{2}\cos 60^\circ\right)=1+\frac{1}{2}=\frac{3}{2}.$ So the area of this triangle is $\frac{\sqrt{3}}{4}s^2=\frac{9\sqrt{3}}{16}.$

 Now let's find the area of the smaller triangles. Notice, $\triangle ACE$ cuts off smaller isosceles triangles from the outer hexagon. The base of these isosceles triangles is perpendicular to the base of the isosceles trapezoid mentioned before, thus we can use trigonometric ratios to find the base and height of these smaller triangles, which are all congruent due to the rotational symmetry of a regular hexagon. The area is then $\frac{1}{2}\left(\frac{1}{2}\cos 60^\circ\right)\left(\frac{1}{2}\sin 60^\circ\right)=\frac{\sqrt{3}}{32}$ and the sum of the areas is $3\cdot \frac{\sqrt{3}}{32}=\frac{3\sqrt{3}}{32}.$

 Therefore, the area of the convex hexagon is $\frac{9\sqrt{3}}{16}-\frac{3\sqrt{3}}{32}=\frac{18\sqrt{3}}{32}-\frac{3\sqrt{3}}{32}=\boxed{\textbf{(C)}\ \frac {15}{32}\sqrt{3}}.$




\textbf{Solution 9 (Linear Transformation)}

If we try to coordinate bash this problem, it's going to look very ugly with a lot of radicals. However, we can alter and skew the diagram in such a way that all ratios of lengths and areas stay the same while making it a lot easier to work with. Then, we can find the ratio of the area of the wanted region to the area of $ABCDEF$ then apply it to the old diagram.
                    \begin{figure}[H]
                        \centering
                        \includegraphics[width=0.5\textwidth]{Solution/AMC_10_10B_2018_24_Solution_Figure_5.png}
                    \end{figure}
                    

                    \begin{figure}[H]
                        \centering
                        \includegraphics[width=0.5\textwidth]{Solution/AMC_10_10B_2018_24_Solution_Figure_6.png}
                    \end{figure}
                    

 The isosceles right triangle with a leg length of $3$ in the new diagram is $\triangle XYZ$ in the old diagram. We see that if we want to take the area of the new hexagon, we must subtract $\frac{3}{4}$ from the area of $\triangle XYZ$ (the red triangles), giving us $\frac{15}{4}.$ However, we need to take the ratio of this area to the area of $ABCDEF,$ which is $\frac{\frac{15}{4}}{12}=\frac{5}{16}.$ Now we know that our answer is $\frac{5}{16} \cdot \frac{3\sqrt{3}}{2}=\boxed{\textbf{(C)}\ \frac {15}{32}\sqrt{3}}.$




\textbf{Solution 10 (Midsegment Theorem)}


                    \begin{figure}[H]
                        \centering
                        \includegraphics[width=0.5\textwidth]{Solution/AMC_10_10B_2018_24_Solution_Figure_7.png}
                    \end{figure}
                    (~MRENTHUSIASM's diagram + modifications)

 \textbf{Lemma 1: $MP$ is midsegment of $\triangle ANE$.} As $ZX$ is a midsegment of trapezoid $FEBA$, we know that $FA||ZX||EB$. Furthermore, by subdividing trapezoid $FEBA$, $ZM,MP,PX$ are midsegments of $\triangle EFA, \triangle ANE, \triangle, ANB$ respectively.

 \textbf{Lemma 2: $[ACE] = \frac12[ABCDEF]$} Due to radial symmetry about the center, $G$, we only have to prove $[AGC] = [CBA]$. As $\triangle AGB$ and $\triangle CBG$ are equilateral triangles, $AC$ therefore divides rhombus $ABCG$ into 2 congruent halves. $\triangle AGC$ and $\triangle CBA$ thus each contain two such congruent triangles, which means that their areas are equal.

 \textbf{Lemma 3: $[AMP] = \frac1{16}[ABCDEF]$} From Lemma 1, $\triangle AMP \simeq \triangle AEN$ with a scale factor of $\frac12$, which means $[AMP] = \frac14[AEN]$. Furthermore, due to symmetry of $\triangle AEN$ and $\triangle CEN$ about $EN$, $[AEN] = \frac12[ACE]$. Therefore, $[AMP] =  \frac18[ACE]$, or (using lemma 2) $\frac1{16}[ABCDEF]$.

 \textbf{Lemma 4: $[ABCDEF] = \frac{3\sqrt3}2$} Because the hexagon is regular, $[ABCDEF] = 6[ABG] = \frac{6\sqrt3}4 = \frac{3\sqrt3}2$. (For more info on where $\frac{\sqrt3}4$ came from, use the search term "area of equilateral triangle")

 \textbf{Final step:} To find the area of the shaded part, we seek $[ACE] - 3[AMP]$. Using lemmas 2 and 3, the expression resolves to $\frac12[ABCDEF] - 3\frac1{16}[ABCDEF] = \frac5{16} \cdot \frac{3\sqrt3}2 = \boxed{\textbf{(C)}\ \frac{15}{32}\sqrt{3}}$.

 


	\item (\textit{AMC 10 10B 2018}) Let $\lfloor x \rfloor$ denote the greatest integer less than or equal to $x$. How many real numbers $x$ satisfy the equation $x^2 + 10,000\lfloor x \rfloor = 10,000x$?


 $\textbf{(A) } 197 \qquad \textbf{(B) } 198 \qquad \textbf{(C) } 199 \qquad \textbf{(D) } 200 \qquad \textbf{(E) } 201$

 \textbf{Answer: }C



\textbf{Solution 1}

This rewrites itself to $x^2=10,000\{x\}$ where $\lfloor x \rfloor + \{x\} = x$.

 Graphing $y=10,000\{x\}$ and $y=x^2$ we see that the former is a set of line segments with slope $10,000$ from $0$ to $1$ with a hole at $x=1$, then $1$ to $2$ with a hole at $x=2$ etc.Here is a graph of $y=x^2$ and $y=16\{x\}$ for visualization.

 
                    \begin{figure}[H]
                        \centering
                        \includegraphics[width=0.5\textwidth]{Solution/AMC_10_10B_2018_25_Solution_Figure_1.png}
                    \end{figure}
                    

 Now notice that when $x=\pm 100$ the graph has a hole at $(\pm 100,10,000)$ which the equation $y=x^2$ passes through and then continues upwards. Thus our set of possible solutions is bounded by $(-100,100)$. We can see that $y=x^2$ intersects each of the lines once and there are $99-(-99)+1=199$ lines for an answer of $\boxed{\text{(C)}~199}$.

 



\textbf{Solution 2}

Same as the first solution, $x^2=10,000\{x\}$.

 
We can write $x$ as $\lfloor x \rfloor+\{x\}$. Expanding everything, we get a quadratic in $\{x\}$ in terms of $\lfloor x \rfloor$:
 \[\{x\}^2+ (2\lfloor x \rfloor -10,000)\{x\} + \lfloor x \rfloor ^2 = 0\]

 
We use the quadratic formula to solve for $\{x\}$ :
 \[\{x\} = \frac {-2\lfloor x \rfloor + 10,000 \pm \sqrt{ ( 2\lfloor x \rfloor - 10,000 )^2 - 4\lfloor x \rfloor^2  }}{2} = \frac {-2\lfloor x \rfloor + 10,000 \pm \sqrt{ 4\lfloor x \rfloor^2 -40,000 \lfloor x \rfloor + 10,000^2- 4\lfloor x \rfloor^2  }}{2}\]

 
Since $0 \leq \{x\} < 1$, we get an inequality which we can then solve. After simplifying a lot, we get that $\lfloor x \rfloor^2 + 2\lfloor x \rfloor - 9999 < 0$.

 
Solving over the integers, $-101 < \lfloor x \rfloor < 99$, and since $\lfloor x \rfloor$ is an integer, there are $\boxed{\text{(C)}~199}$ solutions. Each value of $\lfloor x \rfloor$ should correspond to one value of $x$, so we are done.

 



\textbf{Solution 3 (Fakesolve)}

Let $x = a+k$ where $a$ is the integer part of $x$ and $k$ is the fractional part of $x$.We can then rewrite the problem below:

 $(a+k)^2 + 10000a = 10000(a+k)$

 From here, we get

 $(a+k)^2 + 10000a = 10000a + 10000k$

 Solving for $a+k = x$
$(a+k)^2 = 10000k$
$x = a+k = \pm100\sqrt{k}$

 Because $0 \leq k < 1$, we know that $a+k$ cannot be less than or equal to $-100$ nor greater than or equal to $100$. Therefore:

 $-99 \leq x \leq 99$

 There are $199$ elements in this range, so the answer is $\boxed{\textbf{(C)} \text{ 199}}$.

 
Note (not by author): this solution seems to be invalid at first, because one can not determine whether $x$ is an integer or not. However, it actually works because although $x$ itself might not be an integer, it is very close to one, so there are 199 potential $x$.

 
Another Note (not by author of previous note): we can actually determine that $x$=0 is the only possible integer value of $x$ is we set $x$=$\lfloor x \rfloor$ we end up with $x$=0   ~YJC64002776

 
Yet Another Note EXCEPT THIS ONE'S VERY IMPORTANT, AS IT DISPROVES THE PROOF (also not by the author of any previous notes or of the original solution): this solution claims that $-99 \leq x \leq 99$, which is not entirely true. When $x = 5000-1000\sqrt{26}$, the equation originally stated in the question holds true. $5000-1000\sqrt{26}$ is roughly $-99.0195$ (rounded to the nearest ten thousandth), which is not in the $-99 \leq x \leq 99$ interval. It is probably by mere coincidence that this solution gives the correct answer, as the solution is not rigorous. Additionally, $x$ is not necessarily integer, so it is not claimable that there are 199 elements in the range of $-99 \leq x \leq 99$.

 
One More Note (not by any of the previous authors): The original guy who wrote the solution just wrote the wrong range lol this is a valid solution (at least I'm pretty sure). Given $(a+k)^2 = 10000k$, dividing each side by 10000 gives $\frac{(a+k)^2}{10000} = k$. Our original definition of $k$ states that $0 \leq k < 1$, so $0 \leq \frac{(a+k)^2}{10000} < 1$. This is equivalent to the range $-100 \leq x < 100$ NOT $-99 \leq x < 99$. If you observe the functions behavior, you realize that between every interval of length $1$ there is one solution, except for the intervals that include $0$, where only $0$ is a solution. There are therefore $99+99+1$ solutions in that range which gives the answer. Also $x$ literally doesn't have to be an integer the original constraint was that $x$ is a real number.

 



\textbf{Solution 4}

Notice the given equation is equivalent to $(\lfloor x \rfloor+\{x\})^2=10,000\{x\}$

 Now we know that $\{x\} < 1$ so plugging in $1$ for $\{x\}$ we can find the upper and lower bounds for the values.

 $(\lfloor x \rfloor +1)^2 = 10,000(1)$
$(\lfloor x \rfloor +1) = \pm 100$
$\lfloor x \rfloor = 99, -101$

 And just like $\textbf{Solution 2}$, we see that $-101 < \lfloor x \rfloor < 99$, and since $\lfloor x \rfloor$ is an integer, there are $\boxed{\text{(C)}~199}$ solutions. Each value of $\lfloor x \rfloor$ should correspond to one value of $x$, so we are done.

 
            \begin{itemize}
        \item $\textbf{Note}$ (not by author of any other notes beforehand): this solution is probably the best and most efficient of all. It solves the problems stated in the fakesolves./
            \end{itemize}
            



\textbf{Solution 5}

Firstly, if $x$ is an integer, then $10,000\lfloor x \rfloor=10,000x$, so $x$ must be $0$.

 If $0<x<1$, then we know the following:

 $0<x^2<1$
$10,000\lfloor x \rfloor =0$
$0<10,000x<10,000$

 Therefore, $0<x^2+10,000\lfloor x \rfloor <1$, which overlaps with $0<10,000x<10,000$. This means that there is at least one real solution between $0$ and $1$. Since $x^2+10,000\lfloor x \rfloor$ increases quadratically and $10,000x$ increases linearly, there is only one solution for this case. 

 Similarly, if $1<x<2$, then we know the following:

 $1<x^2<4$
$10,000\lfloor x \rfloor =10,000$
$<10,000<10,000x<20,000$

 By following similar logic, we can find that there is one solution between $1$ and $2$. 

 We can also follow the same process to find that there are negative solutions for $x$ as well.

 There are not an infinite amount of solutions, so at one point there will be no solutions when $n<x<n+1$ for some integer $n$. For there to be no solutions in a given range means that the range of $10,000\lfloor x \rfloor + x^2$ does not intersect the range of $10,000x$. $x^2$ will always be positive, and $10,000\lfloor x \rfloor$ is less than $10,000$ less than $10,000x$, so when $x^2 >= 10,000$, the equation will have no solutions. This means that there are $99$ positive solutions, $99$ negative solutions, and $0$ for a total of $\boxed{\text{(C)}~199}$ solutions.

 



\textbf{Solution 6 (General Equation)}

General solution to this type of equation $f(x, \lfloor x \rfloor) = 0$:

 \begin{verbatim}
1. solve $f(x, \lfloor x \rfloor) = 0$ for $x$ to get $x = g(\lfloor x \rfloor )$2. apply $\lfloor x \rfloor \le x < \lfloor x \rfloor+1$, solve $\lfloor x \rfloor \le g(\lfloor x \rfloor) < \lfloor x \rfloor+1$ to get the domain of $\lfloor x \rfloor$3. get $\lfloor x \rfloor$ from the domain of $\lfloor x \rfloor$ because $\lfloor x \rfloor$ is integer, then get $x$ from $\lfloor x \rfloor$ by $x = g( \lfloor x \rfloor)$Note: function $\lfloor x \rfloor$ maps $x$ to its floor. By solving $f(x, \lfloor x \rfloor) = 0$, we get function $x = g( \lfloor x \rfloor)$, mapping $x$'s floor to $x$
\end{verbatim}

$x^2 - 10000x + 10000 \lfloor x \rfloor =0$
$x=5000 \pm 100 \sqrt{2500- \lfloor x \rfloor}$, $\lfloor x \rfloor \le 2500$
$\lfloor x \rfloor \le x < \lfloor x \rfloor + 1$

 If $x= 5000 + 100 \sqrt{2500 - \lfloor x \rfloor}$, $x \ge 5000$, it contradicts $x < \lfloor x \rfloor + 1 \le 2501$

 So $x= 5000 - 100 \sqrt{2500 - \lfloor x \rfloor}$

 Let $k = \lfloor x \rfloor$ , $x= 5000 - 100 \sqrt{2500 - k}$
$k \le 5000 - 100 \sqrt{2500 - k} < k + 1$
$0 \le 5000 - k - 100 \sqrt{2500 - k} < 1$
$0 \le 2500 - k - 100 \sqrt{2500 - k} + 2500 < 1$
$0 \le (\sqrt{2500 - k} - 50)^2 < 1$
$-1 < \sqrt{2500 - k} - 50 < 1$
$49 < \sqrt{2500 - k} < 51$
$-101 < k < 99$

 So the number of $k$'s values is $99-(-101)-1=199$. Because $x=5000-100\sqrt{2500-k}$, for each value of $k$, there is a value for $x$. The answer is $\boxed{\textbf{(C)} 199}$




\textbf{Solution 7 (Fakesolve)}

Subtracting $10000\lfloor x\rfloor$ from both sides gives $x^2=10000(x-\lfloor x\rfloor)=10000\{x\}$. Dividing both sides by $10000$ gives $\left(\frac{x}{100}\right)^2=\{x\}<1$.  $\left(\frac{x}{100}\right)^2<1$ when $-100<x<100$ so the answer is $\boxed{199}$.

 



\textbf{Solution 8 (Also Gives General Formula For Values of x)}

The question wants to know how many values of $x$ satisfy the equation $x^2 + 10,000\lfloor x \rfloor = 10,000x$. This equation can be simplified as follows

 
 \[x^2 + 10,000\lfloor x \rfloor = 10,000x\]

 \[x^2 = 10,000(x-\lfloor x \rfloor)\]

 Notice that $x-\lfloor x \rfloor$ must be greater than or equal to zero, but less than one. Because of that, the right-hand side of the equation must be less than $10000$. Therefore, $-100<x<100$.

 $x$ can be expressed as $a+b$, where $a$ is an integer and $b$ is a real number such that $0 \leq b < 1$. Notice that in order to satisfy the conditions set down for $x$, $-100 \leq a \leq 99$.

 Substituting $a$ and $b$ into $x^2 = 10000(x-\lfloor x \rfloor)$, we get

 
 \[(a+b)^2 = 10,000b\]

 \[a^2+2ab+b^2 = 10,000b\]

 Let's try to turn this into a quadratic equation where we're trying to solve for $b$. Simplifying, we find

 
 \[b^2+(2a-10,000)b+a^2 = 0\]

 Now the value of $b$ depends on the value of $a$. Our task is to figure out how many values of $a$ will give me valid values of $b$ (a.k.a the values of $a$ that give me $b$ such that $0 \leq b < 1$), as that will be our answer. We must also keep in mind that $-100 \leq a \leq 99$.

 Using the quadratic formula, we can find

 
 \[\frac{(-2a+10,000) \pm \sqrt{(2a-10,000)^2-4a^2}}{2}\]

 The expression above represents the values of $b$ in terms of $a$. Since it represents $b$, we want values of $a$ such that the equation is between zero and one, zero inclusive. We can then conceive the inequality below.

 
 \[0 \leq \frac{(-2a+10,000) \pm \sqrt{(2a-10,000)^2-4a^2}}{2} < 1\]

 Wow, that looks like a mess! How are we supposed to easily use such a complicated expression? Maybe it would help to simplify it further. First, let's notice a few things.

 
            \begin{itemize}
        \item I want the numerator to be between 0 and 2, 0 inclusive. Of the $\pm$ in the quadratic formula, only the "$-$" (minus) will work. Why? Because $-2a+10000$ will always be positive (if I plug in the largest possible value of $a$, $99$, it will give $9802$), and if I want $b$ to be real, which I do, then $\sqrt{(2a-10,000)^2-4a^2}$ should also be positive. $-2a+10000$ will already be much bigger than 2, and adding more on will only make it bigger. It will never be less than 2. It will never work.
            \end{itemize}
            
            \begin{itemize}
        \item I can simplify the numbers under the square root using Difference of Squares. I can turn $\sqrt{(2a-10,000)^2-4a^2}$ into $\sqrt{(2a-10,000)^2-(2a)^2}$. Applying Difference of Squares, it will eventually simplify out to be $200\sqrt{2500-a}$.
            \end{itemize}
            Cool! Now we've simplified the square root and figured out that it must be a minus sign, not a plus sign. Our new inequality looks like this:

 
 \[0 \leq \frac{(-2a+10,000) - 200\sqrt{2500-a}}{2} < 1\]

 Simplifying the unsimplified fraction further, we can get this:

 
 \[0 \leq -a + 5,000 - 100\sqrt{2500-a} < 1\]

 Yay! A simpler equation. As long as the value $a$ fits this inequality, the restrictions set on $b$ are satisfied. But wait. What about the restrictions on $a$?

 \textbf{By testing out $-100$ and $99$ (and finding that $-100$ gives a value very close to $1$, and that $99$ gives exactly $1$), we can conclude that any value of $-100 \leq$a$\leq 99$ will satisfy the inequality above as long as it does not make $-a + 5,000 - 100\sqrt{2500-a}$ equal to $1$. By testing $-100$ and $99$, you will also find that $a = 99$ will not satisfy the inequality.} <u><-- old, non-rigorous way. You can ignore it, but if you want to see a bogus way to prove it, feel free to read it.</u>

 I can solve $0 \leq -a + 5000 - 100\sqrt{2500-a} < 1$ for $a$ this way:

 $0 \leq -a + 5000 - 100\sqrt{2500-a} < 1$
$-5000 \leq -a - 100\sqrt{2500-a} < -4999$
$5000 \geq a + 100\sqrt{2500-a} < 4999$
$5000-a \geq 100\sqrt{2500-a} < 4999-a$
$25000000-10000a+a^2 \geq 25000000-10000a > 4999^2-9998a+a^2$
$a^2 \geq 0 > -9999+2a+a^2$
$0 \geq -a^2 > -9999+2a$
$0 \leq a^2 < 9999-2a$
$2a \leq a(a+2) < 9999$

 Following the left part of the inequality,

 $2a \leq a^2 + 2a$
$0 \leq a^2$
$a$ can be anything $(-\infty, \infty)$.

 Following the right part of the inequality,

 $a(a+2) < 9999$
$a < 99$ and $a > -101$
$a$ can be anything $(-101, 99)$.

 Combining both restrictions and given previously that $-100 \leq a \leq 99$, we can conclude that the inequality describing all possible values of $a$ is $-100 \leq a < 99$.

 Remember how these values of $a$ give valid values of $b$ that satisfy the desired inequality $0 \leq b < 1$ and how adding up $a$ and $b$ gives $x$. In summary, remember that the number of values of $a$ is equivalent to the number of values of $x$.

 $a$ must be integer, so the number of values of $a$ that satisfy $-100 \leq a < 99$ is $\boxed{199}$.

 The answer is $\boxed{\textbf{(C)} 199}$.

 
Note From Author of Solution: you may have noticed that I missed a slight bit by not proving that the thing under the square root in the quadratic formula I made was positive. Don't worry, it wasn't a mistake. It's actually pretty easy to prove. I'll leave it to you to figure out how to do it!

 



\textbf{Solution 9}

Note that the problem rearranges itself to $x^2=10,000\{x\}$ where $\lfloor x \rfloor + \{x\} = x$.

 
To get a general scope of solutions, by noting that $0 \leq \{x\} < 1$, we have that $0 \leq x^2 < 10,000$. Therefore, $-100 < x < 100$.

 
From here, it is easy to observe that there is one solution per each integer value of $-99 \leq \lfloor x \rfloor \leq 99$. This can be reasoned by bounding $\{x\}$ for each $\lfloor x \rfloor$ (intuitively, this is because $\lfloor x \rfloor$ determines the fractional part: for each $\lfloor x \rfloor$, you will be forced to choose the value of $\{x\}$ that will satisfy the equation).

 
Ultimately, the answer is $99+99+1 = \boxed{\textbf{(C)} 199}$.

 


 


\end{enumerate}